\documentclass[12pt]{book} % Change to book
\usepackage{amsmath} % For mathematical expressions
\usepackage{xcolor} % For colored text
\usepackage{graphicx} % For including images
\usepackage{url}
\usepackage{booktabs} % For better-looking tables
\usepackage{makecell}
\usepackage{indentfirst}
\usepackage{algorithm}
\usepackage{algpseudocode}
\usepackage{amssymb}
\usepackage{placeins} % For \FloatBarrier
\usepackage{listings}
% Define the language style For Python code
\lstset{
    language=Python,
    basicstyle=\ttfamily,
    keywordstyle=\color{blue},
    stringstyle=\color{red},
    commentstyle=\color{gray},
    morecomment=[l][\color{magenta}]{\#},
    frame=single,
    breaklines=true
}
\begin{document}

\chapter{Algorithms Analysis}

\section{Introduction to Algorithm Analysis}
In computer science, the \textbf{analysis of algorithms} is the process of finding the computational complexity of algorithms—the amount of time, storage, or other resources needed to execute them. Usually, this involves determining a function that relates the size of an algorithm's input to the number of steps it takes (its time complexity) or the number of storage locations it uses (its space complexity). An algorithm is said to be efficient when this function's values are small or grow slowly compared to a growth in the size of the input. DIfferent inputs of the same size may cause the algorithm to have dIfferent behavior, so best, worst, and average case descriptions might all be of practical interest. When not otherwise specIfied, the function describing the perFormance of an algorithm is usually an upper bound, determined from the worst-case inputs to the algorithm.

 Algorithm analysis is an important part of a broader computational complexity theory, which provides theoretical estimates For the resources needed by any algorithm that solves a given computational problem. These estimates provide insight into reasonable directions of search For efficient algorithms.

In theoretical analysis of algorithms, it is common to estimate their complexity in the asymptotic sense, i.e., to estimate the complexity function For arbitrarily large input. Big O notation, Big-omega (\(\Omega\)) notation, and Big-theta (\(\Theta\)) notation are used to this end. For instance, binary search is said to run in a number of steps proportional to the logarithm of the size $n$ of the sorted list being searched, or in $O(\log n)$, colloquially "in logarithmic time". Usually asymptotic estimates are used because dIfferent implementations of the same algorithm may dIffer in efficiency. However, the efficiencies of any two "reasonable" implementations of a given algorithm are related by a constant multiplicative factor called a 'hidden constant'.

Exact (not asymptotic) measures of efficiency can sometimes be computed but they usually require certain assumptions concerning the particular implementation of the algorithm, called model of computation. A model of computation may be defined in terms of an abstract machine, e.g., Turing machine, and/or by postulating that certain operations are executed in unit time. For example, If the sorted list to which we apply binary search has $n$ elements, and we can guarantee that each lookup of an element in the list can be done in unit time, then at most $\log_2(n) + 1$ time units are needed to Return an answer.

\subsection{Definition and Objectives}
Time efficiency estimates depend on what we define to be a step. For the analysis to correspond usefully to the actual run-time, the time required to perForm a step must be guaranteed to be bounded above by a constant. One must be careful here; For instance, some analyses count an addition of two numbers as one step. This assumption may not be warranted in certain contexts. For example, If the numbers involved in a computation may be arbitrarily large, the time required by a single addition can no longer be assumed to be constant.

Two cost models are generally used~\cite{AhoHopcroft1974, Hromkovič2004, Ausiello1999, Wegener20, Tarjan1983}:
\begin{itemize}
    \item 
    \textbf{UnIForm Cost Model}:
    The unIForm cost model, also called unit-cost model (and similar variations), assigns a constant cost to every machine operation, regardless of the size of the numbers involved. 
    \newline
    Let's consider an example algorithm For finding the sum of elements in an array: 

    \begin{algorithm}[H]
    \caption{Sum of Array Elements}
    \begin{algorithmic}[1]
    \Function{sumArray}{$A, n$}
        \State $sum \gets 0$
        \For{$i \gets 1 \text{ to } n$}
            \State $sum \gets sum + A[i]$
        \EndFor
        \State \Return $sum$
    \EndFunction
    \end{algorithmic}
    \end{algorithm}

    In this algorithm, each iteration of the \texttt{For} loop takes the same amount of time as it perForms a constant number of operations. ThereFore, the unIForm cost model assumes each operation in the loop has unIForm cost.

    \item 
    \textbf{Logarithmic Cost Model}:
    The logarithmic cost model, also called logarithmic-cost measurement (and similar variations), assigns a cost to every machine operation proportional to the number of bits involved. 
    \newline
    Let's consider the example of the binary search algorithm:

    \begin{algorithm}[H]
    \caption{Binary Search}
    \begin{algorithmic}[1]
    \Function{binarySearch}{$A, key, low, high$}
        \If{$low \leq high$}
            \State $mid \gets \frac{low + high}{2}$
            \If{$A[mid] == key$}
                \State \Return $mid$
            \ElsIf{$A[mid] > key$}
                \State \Return \Call{binarySearch}{$A, key, low, mid-1$}
            \Else
                \State \Return \Call{binarySearch}{$A, key, mid+1, high$}
            \EndIf
        \EndIf
        \State \Return $-1$
    \EndFunction
    \end{algorithmic}
    \end{algorithm}

    In binary search, the input size is divided in half at each step, leading to a logarithmic number of steps to find the desired element. ThereFore, the cost of the algorithm is logarithmic in the input size.
\end{itemize}

The latter is more cumbersome to use, so it is only employed when necessary, For example in the analysis of arbitrary-precision arithmetic algorithms, like those used in cryptography.

A key point which is often overlooked is that published lower bounds For problems are often given For a model of computation that is more restricted than the set of operations that you could use in practice and thereFore there are algorithms that are faster than what would naively be thought possible~\cite{PriceAbstraction}.

\subsection{The Role of Algorithm Analysis in Software Development}
Run-time analysis is a theoretical classIfication that estimates and anticipates the increase in running time (or run-time or execution time) of an algorithm as its input size (usually denoted as $n$) increases. Run-time efficiency is a topic of great interest in computer science: A program can take seconds, hours, or even years to finish executing, depending on which algorithm it implements. While software profiling techniques can be used to measure an algorithm's run-time in practice, they cannot provide timing data For all infinitely many possible inputs; the latter can only be achieved by the theoretical methods of run-time analysis.

\subsection{Overview of Complexity Measures}
InFormally, an algorithm can be said to exhibit a growth rate on the order of a mathematical function If beyond a certain input size $n$, the function $f(n)$ times a positive constant provides an upper bound or limit For the run-time of that algorithm. In other words, For a given input size $n$ greater than some $n_0$ and a constant $c$, the run-time of that algorithm will never be larger than $c \times f(n)$. This concept is frequently expressed using Big O notation. For example, since the run-time of insertion sort grows quadratically as its input size increases, insertion sort can be said to be of order $O(n^2)$.

Big O notation is a convenient way to express the worst-case scenario For a given algorithm, although it can also be used to express the average-case -- For example, the worst-case scenario For quicksort is $O(n^2)$, but the average-case run-time is $O(n \log n)$.

Consider the following Insertion Sort related example:\newline
Insertion sort is a simple sorting algorithm that builds the final sorted array one element at a time. It iterates through an array and at each iteration, it removes one element from the input data and finds its correct position within the sorted part of the array.

\begin{algorithm}
\caption{Insertion Sort}
\begin{algorithmic}[1]
\Procedure{InsertionSort}{$A, n$}
    \For{$i \gets 1$ to $n$}
        \State $key \gets A[i]$
        \State $j \gets i - 1$
        \While{$j \geq 0$ and $A[j] > key$}
            \State $A[j + 1] \gets A[j]$
            \State $j \gets j - 1$
        \EndWhile
        \State $A[j + 1] \gets key$
    \EndFor
\EndProcedure
\end{algorithmic}
\end{algorithm}

Let $T(n)$ represent the worst-case run-time of insertion sort For an input of size $n$. In each iteration of the outer loop, insertion sort compares the key with the elements in the sorted subarray and shIfts larger elements one position to the right until the correct position For the key is found. 

Considering the worst-case scenario, when the array is in reverse order, insertion sort would perForm maximum comparisons and shIfts in each iteration. In the worst case, For each $i$ from $2$ to $n$, insertion sort perForms $i-1$ comparisons and shIfts.

Hence, the worst-case run-time $T(n)$ can be expressed as:
$$ T(n) = \sum_{i=2}^{n} (i - 1) = 1 + 2 + 3 + \dots + (n-1) = \frac{n(n-1)}{2} = O(n^2) $$




\subsection{Evaluating run-time complexity}

The run-time complexity For the worst-case scenario of a given algorithm can sometimes be evaluated by examining the structure of the algorithm and making some simplIfying assumptions.  Consider the following pseudocode:

\begin{verbatim}
 1    get a positive integer n from input
 2    If n > 10
 3        print "This might take a While..."
 4    For i = 1 to n
 5        For j = 1 to i
 6            print i * j
 7    print "Done!"
\end{verbatim}

A given computer will take a discrete amount of time to execute each of the instructions involved with carrying out this algorithm.  Say that the actions carried out in step 1 are considered to consume time at most $T_1$, step 2 uses time at most $T_2$, and so Forth.

In the algorithm above, steps 1, 2, and 7 will only be run once.  For a worst-case evaluation, it should be assumed that step 3 will be run as well.  Thus the total amount of time to run steps 1-3 and step 7 is:

\[
T_1 + T_2 + T_3 + T_7.
\]

The loops in steps 4, 5, and 6 are trickier to evaluate.  The outer loop test in step 4 will execute $(n + 1)$ times\footnote{an extra step is required to terminate the For loop, hence $n + 1$ and not $n$ executions}, which will consume $T_4(n + 1)$ time.  The inner loop, on the other hand, is governed by the value of $j$, which iterates from 1 to $i$.  On the first pass through the outer loop, $j$ iterates from 1 to 1:  The inner loop makes one pass, so running the inner loop body (step 6) consumes $T_6$ time, and the inner loop test (step 5) consumes $2T_5$ time.  During the next pass through the outer loop, $j$ iterates from 1 to 2:  the inner loop makes two passes, so running the inner loop body (step 6) consumes $2T_6$ time, and the inner loop test (step 5) consumes $3T_5$ time.

Altogether, the total time required to run the inner loop body can be expressed as an arithmetic progression:

\[
T_6 + 2T_6 + 3T_6 + \cdots + (n-1) T_6 + n T_6
\]

which can be factored\footnote{It can be proven by induction that $1 + 2 + 3 + \cdots + (n-1) + n = \frac{n(n+1)}{2}$} as

\[
\left[ 1 + 2 + 3 + \cdots + (n-1) + n \right] T_6 = \left[ \frac{1}{2} (n^2 + n) \right] T_6
\]

The total time required to run the inner loop \emph{test} can be evaluated similarly:

\begin{align}
& 2T_5 + 3T_5 + 4T_5 + \cdots + (n-1) T_5 + n T_5 + (n + 1) T_5\\
=&\ T_5 + 2T_5 + 3T_5 + 4T_5 + \cdots + (n-1)T_5 + nT_5 + (n+1)T_5 - T_5
\end{align}

which can be factored as

\begin{align}
& T_5 \left[ 1+2+3+\cdots + (n-1) + n + (n + 1) \right] - T_5 \\
=& \left[ \frac{1}{2} (n^2 + n) \right] T_5 + (n + 1)T_5 - T_5 \\
=& \left[ \frac{1}{2} (n^2 + n) \right] T_5 + n T_5 \\
=& \left[ \frac{1}{2} (n^2 + 3n) \right] T_5
\end{align}

ThereFore, the total run-time For this algorithm is:

\[
f(n) = T_1 + T_2 + T_3 + T_7 + (n + 1)T_4 + \left[ \frac{1}{2} (n^2 + n) \right] T_6 + \left[ \frac{1}{2} (n^2+3n) \right] T_5
\]

which reduces to

\[
f(n) = \left[ \frac{1}{2} (n^2 + n) \right] T_6 + \left[ \frac{1}{2} (n^2 + 3n) \right] T_5 + (n + 1)T_4 + T_1 + T_2 + T_3 + T_7
\]

As a rule-of-thumb, one can assume that the highest-order term in any given function dominates its rate of growth and thus defines its run-time order.  In this example, $n^2$ is the highest-order term, so one can conclude that $f(n) = O(n^2)$.  Formally this can be proven as follows:

\begin{quote}
Prove that $\left[ \frac{1}{2} (n^2 + n) \right] T_6 + \left[ \frac{1}{2} (n^2 + 3n) \right] T_5 + (n + 1)T_4 + T_1 + T_2 + T_3 + T_7 \le cn^2,\ n \ge n_0$

\begin{align}
%\[
&\left[ \frac{1}{2} (n^2 + n) \right] T_6 + \left[ \frac{1}{2} (n^2 + 3n) \right] T_5 + (n + 1)T_4 + T_1 + T_2 + T_3 + T_7\\
\le &( n^2 + n )T_6 + ( n^2 + 3n )T_5 + (n + 1)T_4 + T_1 + T_2 + T_3 + T_7 \ (\text{For } n \ge 0 )
%\]
\end{align}


Let $k$ be a constant greater than or equal to $[T_1..T_7]$

\begin{align}
&T_6( n^2 + n ) + T_5( n^2 + 3n ) + (n + 1)T_4 + T_1 + T_2 + T_3 + T_7 \le k( n^2 + n ) + k( n^2 + 3n ) + kn + 5k\\
= &2kn^2 + 5kn + 5k \le 2kn^2 + 5kn^2 + 5kn^2 + 5kn \le 12kn^2
\end{align}
For $n \ge 1$.

ThereFore, $f(n) = O(n^2)$.
\end{quote}

\section{Understanding Complexity Notations}
In theoretical analysis of algorithms, it is common to estimate their complexity in the asymptotic sense, i.e., to estimate the complexity function For arbitrarily large input. Big O notation, Big Omega (\(\Omega\)) notation, and Big Theta (\(\Theta\)) notation are used to this end. For instance, binary search is said to run in a number of steps proportional to the logarithm of the size $n$ of the sorted list being searched, or in $O(\log n)$, colloquially 'in logarithmic time'. Usually asymptotic estimates are used because dIfferent implementations of the same algorithm may dIffer in efficiency. However, the efficiencies of any two 'reasonable' implementations of a given algorithm are related by a constant multiplicative factor called a 'hidden constant'.

\subsection{The Big O Notation}
\subsubsection{Definition and SignIficance}

Big O notation is a mathematical notation that describes the limiting behavior of a function when the argument tends towards a particular value or infinity. In the context of algorithm analysis, Big O notation is used to describe the upper bound on the growth rate of an algorithm's time or space complexity. It provides a standardized and concise way to express the efficiency of algorithms by focusing on their worst-case perFormance.

The signIficance of Big O notation lies in its ability to provide a clear understanding of how an algorithm scales with input size. By analyzing the upper bound of an algorithm's complexity, engineers and developers can make inFormed decisions about algorithm selection and optimization strategies. Additionally, Big O notation facilitates communication and comparison between dIfferent algorithms, enabling the identIfication of the most efficient solutions For specIfic problem domains.

\subsubsection{Basic Examples}

To illustrate the concept of Big O notation, consider the following basic examples:

\begin{itemize}
    \item \textbf{Constant Time}: An algorithm that perForms a single operation regardless of the input size has a time complexity of O(1). For example, accessing an element in an array by index or checking If a number is even or odd.
    \FloatBarrier
    Consider the following algorithm:
    \FloatBarrier
    \begin{algorithm}
\caption{Access Element in Array}
\begin{algorithmic}[1]
\Function{AccessElement}{$arr, index$}
    \State \textbf{Return} $arr[index]$
\EndFunction
\end{algorithmic}
\end{algorithm}
\FloatBarrier
    \item \textbf{Linear Time}: An algorithm whose runtime grows linearly with the input size has a time complexity of O(n). For example, iterating through each element in an array or list.
    \FloatBarrier
    Following is the algorithmic example:
\FloatBarrier
\begin{algorithm}[H]
\caption{Example Algorithm For Big O Notation}
\begin{algorithmic}
\State function foo($n$)
\State $sum = 0$
\For{$i = 1$ to $n$}
\State $sum = sum + i$
\EndFor
\State Return $sum$
\end{algorithmic}
\end{algorithm}
\FloatBarrier
The time complexity of this algorithm is $O(n)$ because the loop runs $n$ times, resulting in a linear growth rate with respect to the input size $n$.
    \FloatBarrier
    \item \textbf{Quadratic Time}: An algorithm whose runtime grows quadratically with the input size has a time complexity of O($n^2$). For example, nested loops where each iteration depends on the size of the input.
    \FloatBarrier
    Following is the algorithmic example:
    \FloatBarrier
    \begin{algorithm}
\caption{Nested Loop}
\begin{algorithmic}[1]
\Function{NestedLoop}{$n$}
    \For{$i \gets 1$ to $n$}
        \For{$j \gets 1$ to $n$}
            \State PerForm some operation
        \EndFor
    \EndFor
\EndFunction
\end{algorithmic}
\end{algorithm}
\FloatBarrier
\end{itemize}






\subsection{The Big Omega (\(\Omega\)) Notation}

\subsubsection{Definition and SignIficance}

Big Omega notation is another mathematical notation used in algorithm analysis. It describes the lower bound on the growth rate of an algorithm's time or space complexity. Essentially, Big Omega notation provides inFormation about the best-case scenario For the algorithm's perFormance. Unlike Big O notation, which represents the upper bound on the growth rate, Big Omega notation focuses on the lower bound, indicating the minimum amount of resources required by the algorithm.

The signIficance of Big Omega notation lies in its ability to provide insights into the inherent efficiency of an algorithm. By analyzing the lower bound of an algorithm's complexity, engineers and developers can gain a better understanding of its best-case perFormance and identIfy situations where the algorithm perForms optimally. Additionally, Big Omega notation complements Big O notation by offering a more comprehensive view of an algorithm's behavior across dIfferent input sizes.

\subsubsection{Basic Examples}

To illustrate the concept of Big Omega notation, consider the following basic examples:

\begin{itemize}
    \item \textbf{Linear Time}: An algorithm whose runtime grows linearly with the input size has a lower bound time complexity of $\Omega(n)$. For example, a linear search algorithm always requires at least one comparison For each element in the input array.
    \FloatBarrier
Consider the following algorithm:
\FloatBarrier
\begin{algorithm}[H]
\caption{Example Algorithm For Big Omega (\(\Omega\)) Notation}
\begin{algorithmic}
\State function bar($n$)
\State $min = A[1]$
\For{$i = 2$ to $n$}
\If{$A[i] < min$}
\State $min = A[i]$
\EndIf
\EndFor
\State Return $min$
\end{algorithmic}
\end{algorithm}
\FloatBarrier
The time complexity of this algorithm is $\Omega(n)$ because in the best-case scenario, the algorithm needs to iterate through the array once to find the minimum value, resulting in a linear growth rate.
    
    \item \textbf{Quadratic Time}: An algorithm whose runtime grows quadratically with the input size has a lower bound time complexity of $\Omega(n^2)$. For example, an algorithm that computes all pairs of elements in an input array requires at least quadratic time to complete.
    \FloatBarrier
    Following is the algorithmic example:
    \FloatBarrier
    \begin{algorithm}
\caption{Quadratic Time Algorithm}
\begin{algorithmic}[1]
\Function{QuadraticTime}{$A, n$}
    \State $result \gets 0$
    \For{$i \gets 1$ to $n$}
        \For{$j \gets 1$ to $n$}
            \State $result \gets result + A[i] \times A[j]$
        \EndFor
    \EndFor
    \State \textbf{Return} $result$
\EndFunction
\end{algorithmic}
\end{algorithm}
\end{itemize}

\subsection{The Big Theta (\(\Theta\)) Notation}

\subsubsection{Definition and SignIficance}

Big Theta (\(\Theta\)) notation is a mathematical notation used in algorithm analysis. It provides both the upper and lower bounds on the growth rate of an algorithm's time or space complexity. Unlike Big O and Big Omega notations, which represent only the upper and lower bounds respectively, Big Theta notation represents the tightest bound on the algorithm's perFormance. This notation is often used when the upper and lower bounds are equal, indicating that the algorithm's perFormance behaves consistently across dIfferent input sizes.

The signIficance of Big Theta notation lies in its ability to precisely describe the behavior of an algorithm. By providing both upper and lower bounds, Big Theta notation offers a clear understanding of the algorithm's perFormance characteristics and helps in determining its efficiency in the worst-case, best-case, and average-case scenarios.

\subsubsection{Understanding the tight bound}

In Big Theta (\(\Theta\)) notation, the upper and lower bounds on an algorithm's perFormance are equal, indicating a tight bound on its growth rate. This implies that the algorithm's time or space complexity grows at the same rate regardless of the input size. In other words, the algorithm's behavior is consistent and predictable, making Big Theta notation particularly useful For analyzing algorithms with well-defined perFormance characteristics.


\subsubsection{Basic Examples}

To illustrate the concept of Big Theta notation, consider the following basic examples:

\begin{itemize}
    \item \textbf{Linear Time}: An algorithm whose runtime grows linearly with the input size has a time complexity of \(\Theta(n)\). For example, a linear search algorithm typically perForms \(n\) comparisons in the worst-case scenario, where \(n\) is the size of the input array.
    \FloatBarrier
    Consider the following algorithm:
\FloatBarrier
\begin{algorithm}[H]
\caption{Example Algorithm For Big Theta (\(\Theta\)) Notation}
\begin{algorithmic}
\State function baz($n$)
\State $sum = 0$
\For{$i = 1$ to $n$}
\State $sum = sum + i$
\EndFor
\For{$j = 1$ to $n$}
\State $sum = sum + j$
\EndFor
\State Return $sum$
\end{algorithmic}
\end{algorithm}
\FloatBarrier
The time complexity of this algorithm is $\Theta(n)$ because it consists of two loops, each running $n$ times. As a result, the algorithm has a linear growth rate that matches both the lower and upper bounds.
\FloatBarrier
    \item \textbf{Quadratic Time}: An algorithm whose runtime grows quadratically with the input size has a time complexity of \(\Theta(n^2)\). For example, an algorithm that computes all pairs of elements in an input array perForms \(\frac{n \times (n-1)}{2}\) operations, where \(n\) is the size of the input array.
    \FloatBarrier
    Following is the algorithmic example:
    \FloatBarrier
    \begin{algorithm}
\caption{Quadratic Time Algorithm}
\begin{algorithmic}[1]
\Function{QuadraticTime}{$A, n$}
    \State $result \gets 0$
    \For{$i \gets 1$ to $n$}
        \For{$j \gets 1$ to $n$}
            \State $result \gets result + A[i] \times A[j]$
        \EndFor
    \EndFor
    \State \textbf{Return} $result$
\EndFunction
\end{algorithmic}
\end{algorithm}
\end{itemize}

\subsection{Comparison of the Three Notations}

When comparing Big O, Big Omega (\(\Omega\)), and Big Theta (\(\Theta\)) notations, it is important to understand their relationships. Big O notation provides an upper bound on the growth rate of a function, Big Omega (\(\Omega\))notation provides a lower bound, and Big Theta (\(\Theta\)) notation provides a tight bound. In many cases, these notations are used together to provide a comprehensive analysis of an algorithm's perFormance.


In summary:
\begin{itemize}
    \item Big O notation: Upper bound
    \item Big Omega (\(\Omega\)) notation: Lower bound
    \item Big Theta (\(\Theta\)) notation: Tight bound
\end{itemize}

\section{Analytical Tools For Algorithm Analysis}

Analytical tools play a crucial role in analyzing the perFormance and efficiency of algorithms. They provide us with a quantIfiable way to understand how algorithms behave under dIfferent scenarios. In this section, we will discuss some fundamental analytical tools For algorithm analysis.

\subsection{Asymptotic Analysis}

Asymptotic analysis is a method used to evaluate the perFormance of an algorithm as the input size tends towards infinity. It helps us understand how the algorithm will scale with larger inputs.

\subsubsection{Growth of Functions}

In asymptotic analysis, we focus on the growth rates of functions that represent the running time or space complexity of algorithms. Common growth rates include constant time, logarithmic time, linear time, quadratic time, etc.

\subsubsection{Common Asymptotic Functions}

Some common asymptotic functions used in algorithm analysis include:
\begin{itemize}
  \item $O(\cdot)$ - Big O notation For the upper bound of a function.
  \item $\Omega(\cdot)$ - Omega notation For the lower bound of a function.
  \item $\Theta(\cdot)$ - Theta notation For tight bounds of a function.
\end{itemize}

\subsection{Recurrence Relations}

Recurrence relations are equations that describe the runtime of a recursive algorithm in terms of its subproblem sizes. Solving these recurrences helps us analyze the time complexity of recursive algorithms.

\subsubsection{Solving Recurrences}

There are various methods to solve recurrence relations, such as substitution method, master theorem, recursion tree method, etc. explained as follows:
\FloatBarrier
Substitution Method:
\FloatBarrier
The substitution method is a technique used to solve recurrence relations by making an educated guess and then proving the correctness of the guess using mathematical induction. Here's how it works:

\begin{enumerate}
    \item \textbf{Guess a solution}: We make an educated guess about the Form of the solution to the recurrence relation. This guess often involves an expression involving the size of the input and some unknown constants.
    
    \item \textbf{Prove by induction}: We then prove the correctness of our guess using mathematical induction. This involves verIfying that the guessed solution satisfies the original recurrence relation and its initial conditions.
    
    \item \textbf{Example}:
    
    Let's consider the recurrence relation For the Fibonacci sequence:
    
    \[
    T(n) = T(n-1) + T(n-2)
    \]
    
    where \(T(n)\) represents the time complexity of computing the \(n\)-th Fibonacci number. By guessing that \(T(n) = O(2^n)\), we can prove by induction that this guess is correct.
\end{enumerate}
\FloatBarrier
Master Theorem:
\FloatBarrier

The master theorem is a powerful tool For solving recurrence relations that arise in the analysis of divide-and-conquer algorithms. It provides a straightForward way to determine the asymptotic behavior of the solution without the need For guesswork or induction. The master theorem applies to recurrence relations of the Form \(T(n) = aT(n/b) + f(n)\), where \(a \geq 1\), \(b > 1\), and \(f(n)\) is a given function.

\begin{enumerate}
    \item \textbf{IdentIfy the Form}: We start by identIfying the Form of the recurrence relation and determining the values of \(a\), \(b\), and \(f(n)\).
    
    \item \textbf{Apply the master theorem}: Based on the values of \(a\), \(b\), and the Form of \(f(n)\), we can apply one of the cases of the master theorem to determine the asymptotic behavior of \(T(n)\).
    
    \item \textbf{Example}:
    
    Consider the recurrence relation \(T(n) = 2T(n/2) + n\). By identIfying \(a = 2\), \(b = 2\), and \(f(n) = n\), we can apply case 2 of the master theorem, which tells us that \(T(n) = \Theta(n \log n)\).
\end{enumerate}
\FloatBarrier
Recursion Tree Method:
\FloatBarrier

The recursion tree method is a graphical technique used to visualize and solve recurrence relations. It involves drawing a tree diagram to represent the recursive calls made by the algorithm and then analyzing the structure of the tree to determine the overall time complexity.

\begin{enumerate}
    \item \textbf{Construct the recursion tree}: We start by drawing a tree diagram representing the recursive calls made by the algorithm. Each node in the tree corresponds to a subproblem, and the edges represent the recursive calls.
    
    \item \textbf{Analyze the tree}: We then analyze the structure of the tree to determine the total number of nodes and levels. By summing up the work done at each level of the tree, we can derive the overall time complexity of the algorithm.
    
    \item \textbf{Example}:
    
    Let's use the recursion tree method to analyze the time complexity of the merge sort algorithm. By constructing a recursion tree For merge sort, we can see that the height of the tree is \(\log n\), and each level of the tree requires \(O(n)\) work. ThereFore, the overall time complexity of merge sort is \(O(n \log n)\).
\end{enumerate}

\subsection{Amortized Analysis}

Amortized analysis is a technique used to determine the average time complexity of a sequence of operations in data structures where some operations may be expensive but are offset by others that are relatively cheap.

\subsubsection{Definition and Applications}

Amortized analysis is a technique used to determine the average time complexity of a sequence of operations in data structures where some operations may be expensive but are offset by others that are relatively cheap. It provides a way to analyze the overall perFormance of data structures by averaging the cost of individual operations over a series of operations.
\FloatBarrier
\textbf{Conceptual Understanding:}
\FloatBarrier

The basic idea behind amortized analysis is to distribute the cost of expensive operations over multiple cheaper operations. This ensures that the average cost of operations remains low even If some individual operations are expensive. 

Consider a dynamic array that doubles its size whenever it runs out of space. Although resizing the array is an expensive operation, it occurs infrequently compared to the cheap constant-time operations of accessing or appending elements. Amortized analysis allows us to show that the average cost of resizing the array is low when spread across multiple append operations.

\FloatBarrier
\textbf{Algorithmic Example: Dynamic Array}
\FloatBarrier

Let's consider a dynamic array that doubles its size whenever it runs out of space. We'll use amortized analysis to analyze the average time complexity of appending \(n\) elements to the array.
\FloatBarrier
\begin{algorithm}
\caption{Append Operation on Dynamic Array}
\begin{algorithmic}[1]
\Function{Append}{$\text{array}, \text{element}$}
    \If{$\text{size}(\text{array}) = \text{capacity}(\text{array})$}
        \State \text{resize}($\text{array}, 2 \times \text{capacity}(\text{array})$)
    \EndIf
    \State $\text{array}[\text{size}(\text{array})] \gets \text{element}$
    \State $\text{size}(\text{array}) \gets \text{size}(\text{array}) + 1$
\EndFunction
\end{algorithmic}
\end{algorithm}
\FloatBarrier
In this example, resizing the array takes \(O(n)\) time, where \(n\) is the current size of the array. However, we can show that the amortized cost of each append operation is \(O(1)\) using the potential method.

\FloatBarrier
\textbf{Applications:}
\FloatBarrier
Amortized analysis is commonly used in the analysis of data structures and algorithms where the cost of individual operations varies signIficantly. Some common applications include:

\begin{itemize}
    \item \textbf{Dynamic arrays}: As demonstrated in the example above, amortized analysis helps in analyzing the average time complexity of dynamic array operations like appending elements.
    
    \item \textbf{Binary heaps}: Amortized analysis can be used to analyze the time complexity of heap operations such as insertion and deletion, where some operations may require restructuring the heap to maintain its properties.
    
    \item \textbf{Hash tables}: In hash tables, resizing the underlying array can be an expensive operation. Amortized analysis helps in analyzing the average time complexity of insertion and deletion operations, taking into account occasional resizing.
\end{itemize}

Overall, amortized analysis provides a useful tool For understanding the average perFormance of algorithms and data structures over a sequence of operations.

\section{Practical Analysis of Algorithms}
Analyzing algorithms is crucial in understanding their efficiency and perFormance. Practical analysis of algorithms involves studying their behavior, runtime complexity, and space complexity using mathematical and empirical techniques.
\subsection{Binary Search}
Binary search is a fundamental searching algorithm that efficiently locates a target value within a sorted array. It works by repeatedly dividing the search interval in half. The key to its efficiency lies in discarding half of the search space at each step.
\subsubsection{Algorithmic Description:}
\begin{algorithm}
\caption{Binary Search}
\begin{algorithmic}[1]
\Function{BinarySearch}{$A[], n, x$}
\State $low \gets 0$
\State $high \gets n - 1$
\While{$low \leq high$}
    \State $mid \gets \lfloor(low + high) / 2 \rfloor$
    \If{$A[mid] == x$}
        \State \Return $mid$
    \ElsIf{$A[mid] < x$}
        \State $low \gets mid + 1$
    \Else
        \State $high \gets mid - 1$
    \EndIf
\EndWhile
\State \Return $-1$ \Comment{Element not found}
\EndFunction
\end{algorithmic}
\end{algorithm}

\subsubsection{Python Implementation of the Binary Search Algorithm}
Below is the Python implementation of the binary search algorithm:

\begin{lstlisting}[language=Python]
def binary_search(arr, target):
    left, right = 0, len(arr) - 1
    
    While left <= right:
        mid = left + (right - left) // 2
        
        If arr[mid] == target:
            Return mid
        elIf arr[mid] < target:
            left = mid + 1
        else:
            right = mid - 1
    
    Return -1
\end{lstlisting}

This function takes a sorted array \texttt{arr} and a target value \texttt{target} as input and Returns the index of \texttt{target} in the array If found, otherwise Returns -1.
\subsubsection{Analysis and Time Complexity}
\textbf{Time Complexity Analysis}

To analyze the time complexity of binary search, let's denote the length of the input array as \(n\). In each iteration of the binary search algorithm, the search interval is divided in half. ThereFore, the size of the search interval is reduced by a factor of 2 in each iteration.

Let \(k\) be the number of iterations required to find the target value or determine that it is not present in the array. Since the search interval size is halved in each iteration, we have:

\[
\frac{n}{2^k} = 1
\]

Solving For \(k\), we get \(k = \log_2 n\). Thus, the time complexity of binary search is \(O(\log n)\).

\textbf{Space Complexity Analysis}

Binary search is an in-place algorithm that does not require any additional memory proportional to the input size \(n\). ThereFore, the space complexity of binary search is \(O(1)\).

Binary search is a highly efficient algorithm For finding elements in a sorted array, with a time complexity of \(O(\log n)\), making it suitable For large datasets where efficiency is crucial.

\subsection{Sorting Algorithms}
Sorting algorithms are crucial in arranging elements in a specIfic order. They are used in various applications like databases, search algorithms, and more. DIfferent sorting algorithms have dIfferent time complexities, influencing their efficiency.

There are numerous sorting algorithms, each with its own advantages and disadvantages. Some of the commonly used sorting algorithms include:

\begin{itemize}
    \item \textbf{Bubble Sort}: Bubble sort repeatedly steps through the list, compares adjacent elements, and swaps them If they are in the wrong order. It continues until the list is sorted.

     Algorithmic Example For Bubble Sort:
     \FloatBarrier % Ensure algorithm appears after the paragraph
    \begin{algorithm}
\caption{Bubble Sort}
\begin{algorithmic}[1]
\Function{BubbleSort}{$A[]$}
    \State $n \gets$ length of $A$
    \For{$i \gets 0$ \textbf{to} $n-1$}
        \For{$j \gets 0$ \textbf{to} $n-i-1$}
            \If{$A[j] > A[j+1]$}
                \State Swap $A[j]$ and $A[j+1]$
            \EndIf
        \EndFor
    \EndFor
\EndFunction
\end{algorithmic}
\end{algorithm}
\FloatBarrier % Ensure algorithm appears after the paragraph
Python Code:
\FloatBarrier % Ensure algorithm appears after the paragraph
\begin{verbatim}
def bubble_sort(arr):
    n = len(arr)
    For i in range(n - 1):
        For j in range(n - i - 1):
            If arr[j] > arr[j + 1]:
                arr[j], arr[j + 1] = arr[j + 1], arr[j]
\end{verbatim}
    \item \textbf{Selection Sort}: Selection sort divides the input list into two parts: the sorted sublist and the unsorted sublist. It repeatedly selects the minimum element from the unsorted sublist and moves it to the beginning of the sorted sublist.
    
    Algorithmic Example For Selection Sort:
     \FloatBarrier % Ensure algorithm appears after the paragraph
     \begin{algorithm}
\caption{Selection Sort}
\begin{algorithmic}[1]
\Function{SelectionSort}{$A[]$}
    \State $n \gets$ length of $A$
    \For{$i \gets 0$ \textbf{to} $n-1$}
        \State $min\_idx \gets i$
        \For{$j \gets i+1$ \textbf{to} $n$}
            \If{$A[j] < A[min\_idx]$}
                \State $min\_idx \gets j$
            \EndIf
        \EndFor
        \State Swap $A[min\_idx]$ and $A[i]$
    \EndFor
\EndFunction
\end{algorithmic}
\end{algorithm}
\FloatBarrier % Ensure algorithm appears after the paragraph
Python Code:
\begin{verbatim}
def selection_sort(arr):
    n = len(arr)
    For i in range(n - 1):
        min_idx = i
        For j in range(i + 1, n):
            If arr[j] < arr[min_idx]:
                min_idx = j
        arr[i], arr[min_idx] = arr[min_idx], arr[i]
\end{verbatim}
\FloatBarrier % Ensure algorithm appears after the paragraph
    
    \item \textbf{Insertion Sort}: Insertion sort builds the final sorted array one item at a time by repeatedly taking the next element from the unsorted part and inserting it into its correct position in the sorted part.
    
    Algorithmic Example For Insertion Sort:
     \FloatBarrier % Ensure algorithm appears after the paragraph
     \begin{algorithm}
\caption{Insertion Sort}
     \begin{algorithmic}[1]
\Function{Insertion Sort}{$A[]$}
    \State $n \gets$ length of $A$
    \For{$i \gets 1$ \textbf{to} $n-1$}
        \State $key \gets A[i]$
        \State $j \gets i - 1$
        \While{$j \geq 0$ and $A[j] > key$}
            \State $A[j + 1] \gets A[j]$
            \State $j \gets j - 1$
        \EndWhile
        \State $A[j + 1] \gets key$
    \EndFor
\EndFunction
\end{algorithmic}
\end{algorithm}
\FloatBarrier % Ensure algorithm appears after the paragraph
Python Code:
\begin{verbatim}
def insertion_sort(arr):
    n = len(arr)
    For i in range(1, n):
        key = arr[i]
        j = i - 1
        While j >= 0 and arr[j] > key:
            arr[j + 1] = arr[j]
            j -= 1
        arr[j + 1] = key
\end{verbatim}
\FloatBarrier % Ensure algorithm appears after the paragraph
    
    \item \textbf{Merge Sort}: Merge sort follows the divide-and-conquer strategy. It divides the input array into two halves, sorts the two halves separately, and then merges them.

 Algorithmic Example For Merge Sort:
     \FloatBarrier % Ensure algorithm appears after the paragraph
\begin{algorithm}
\caption{Merge Sort}
\begin{algorithmic}[1]
\Function{MergeSort}{$A[], start, end$}
    \If{$start < end$}
        \State $mid \gets (start + end) // 2$
        \State \Call{MergeSort}{$A[], start, mid$}
        \State \Call{MergeSort}{$A[], mid+1, end$}
        \State \Call{Merge}{$A[], start, mid, end$}
    \EndIf
\EndFunction
\FloatBarrier % Ensure algorithm appears after the paragraph
\Function{Merge}{$A[], start, mid, end$}
    \State $n_1 \gets mid - start + 1$
    \State $n_2 \gets end - mid$
    \State Create temporary arrays $L[1 \dots n_1]$ and $R[1 \dots n_2]$
    \For{$i \gets 1$ \textbf{to} $n_1$}
        \State $L[i] \gets A[start + i - 1]$
    \EndFor
    \For{$j \gets 1$ \textbf{to} $n_2$}
        \State $R[j] \gets A[mid + j]$
    \EndFor
    \State $i \gets 1$, $j \gets 1$, $k \gets start$
    \While{$i \leq n_1$ and $j \leq n_2$}
        \If{$L[i] \leq R[j]$}
            \State $A[k] \gets L[i]$
            \State $i \gets i + 1$
        \Else
            \State $A[k] \gets R[j]$
            \State $j \gets j + 1$
        \EndIf
        \State $k \gets k + 1$
    \EndWhile
    \While{$i \leq n_1$}
        \State $A[k] \gets L[i]$
        \State $i \gets i + 1$, $k \gets k + 1$
    \EndWhile
    \While{$j \leq n_2$}
        \State $A[k] \gets R[j]$
        \State $j \gets j + 1$, $k \gets k + 1$
    \EndWhile
\EndFunction
\end{algorithmic}
\end{algorithm}
\FloatBarrier % Ensure algorithm appears after the paragraph
Python Code:
\FloatBarrier % Ensure algorithm appears after the paragraph
\begin{verbatim}
def merge_sort(arr):
    If len(arr) > 1:
        mid = len(arr) // 2
        L = arr[:mid]
        R = arr[mid:]

        merge_sort(L)
        merge_sort(R)

        i = j = k = 0

        While i < len(L) and j < len(R):
            If L[i] < R[j]:
                arr[k] = L[i]
                i += 1
            else:
                arr[k] = R[j]
                j += 1
            k += 1

        While i < len(L):
            arr[k] = L[i]
            i += 1
            k += 1

        While j < len(R):
            arr[k] = R[j]
            j += 1
            k += 1
\end{verbatim}
    \FloatBarrier % Ensure algorithm appears after the paragraph
    \item \textbf{Quick Sort}: Quick sort also follows the divide-and-conquer strategy. It selects a pivot element and partitions the array into two subarrays such that elements less than the pivot are on the left and elements greater than the pivot are on the right. It then recursively sorts the subarrays.

Algorithmic Example For Quick Sort:
     \FloatBarrier % Ensure algorithm appears after the paragraph
    \begin{algorithm}
\caption{Quick Sort}
\begin{algorithmic}[1]
\Function{QuickSort}{$A[], start, end$}
    \If{$start < end$}
        \State $pivot \gets$ Partition($A[], start, end$)
        \State \Call{QuickSort}{$A[], start, pivot-1$}
        \State \Call{QuickSort}{$A[], pivot+1, end$}
    \EndIf
\EndFunction
\FloatBarrier % Ensure algorithm appears after the paragraph
\Function{Partition}{$A[], start, end$}
    \State $pivot \gets A[end]$
    \State $i \gets start - 1$
    \For{$j \gets start$ \textbf{to} $end-1$}
        \If{$A[j] < pivot$}
            \State $i \gets i + 1$
            \State Swap $A[i]$ and $A[j]$
        \EndIf
    \EndFor
    \State Swap $A[i + 1]$ and $A[end]$
    \State \Return $i + 1$
\EndFunction
\end{algorithmic}
\end{algorithm}
\FloatBarrier % Ensure algorithm appears after the paragraph
Python Code:
\begin{verbatim}
def quick_sort(arr, start, end):
    If start < end:
        pivot = partition(arr, start, end)
        quick_sort(arr, start, pivot - 1)
        quick_sort(arr, pivot + 1, end)

def partition(arr, start, end):
    pivot = arr[end]
    i = start - 1
    For j in range(start, end):
        If arr[j] < pivot:
            i += 1
            arr[i], arr[j] = arr[j], arr[i]
    arr[i + 1], arr[end] = arr[end], arr[i + 1]
    Return i + 1
\end{verbatim}
    \FloatBarrier % Ensure algorithm appears after the paragraph
    \item \textbf{Heap Sort}: Heap sort builds a heap data structure from the input array and repeatedly extracts the maximum element from the heap and rebuilds the heap until the array is sorted.
    Algorithmic Example For Heap Sort:
     \FloatBarrier % Ensure algorithm appears after the paragraph
     \begin{algorithm}
\caption{Heap Sort}
\begin{algorithmic}[1]
\Function{HeapSort}{$A[]$}
    \State BuildMaxHeap($A[]$)
    \For{$i \gets$ length of $A$ down to $2$}
        \State Swap $A[1]$ and $A[i]$
        \State HeapIfy($A[], 1, i-1$)
    \EndFor
\EndFunction
\FloatBarrier % Ensure algorithm appears after the paragraph
\Function{BuildMaxHeap}{$A[]$}
    \State $n \gets$ length of $A$
    \For{$i \gets n // 2$ down to $1$}
        \State HeapIfy($A[], i, n$)
    \EndFor
\EndFunction
\FloatBarrier % Ensure algorithm appears after the paragraph
\Function{HeapIfy}{$A[], i, n$}
    \State $largest \gets i$
    \State $l \gets 2i$
    \State $r \gets 2i + 1$
    \If{$l \leq n$ and $A[l] > A[largest]$}
        \State $largest \gets l$
    \EndIf
    \If{$r \leq n$ and $A[r] > A[largest]$}
        \State $largest \gets r$
    \EndIf
    \If{$largest \neq i$}
        \State Swap $A[i]$ and $A[largest]$
        \State HeapIfy($A[], largest, n$)
    \EndIf
\EndFunction
\end{algorithmic}
\end{algorithm}
\FloatBarrier % Ensure algorithm appears after the paragraph
Python Code:
\FloatBarrier % Ensure algorithm appears after the paragraph
\begin{verbatim}
def heap_sort(arr):
    n = len(arr)
    For i in range(n // 2 - 1, -1, -1):
        heapIfy(arr, n, i)
    For i in range(n - 1, 0, -1):
        arr[i], arr[0] = arr[0], arr[i]
        heapIfy(arr, i, 0)

def heapIfy(arr, n, i):
    largest = i
    l = 2 * i + 1
    r = 2 * i + 2
    If l < n and arr[l] > arr[largest]:
        largest = l
    If r < n and arr[r] > arr[largest]:
        largest = r
    If largest != i:
        arr[i], arr[largest] = arr[largest], arr[i]
        heapIfy(arr, n, largest)
\end{verbatim}
\end{itemize}
\FloatBarrier % Ensure algorithm appears after the paragraph
\subsubsection{Complexity Analysis of the sorting algorithms}

The space and time complexity of sorting algorithms varies depending on the algorithm used and the input data. Here are the average and worst-case time complexities along wiht the space complexities of some common sorting algorithms:

\begin{itemize}
    \item \textbf{Bubble Sort}: 
    \begin{itemize}
        \item \textbf{Time Complexity: }$O(n^2)$
        \item \textbf{Space Complexity: }Bubble sort is an in-place comparison-based sorting algorithm. It only requires a constant amount of additional space For storing temporary variables. ThereFore, the space complexity of bubble sort is \(O(1)\).
    \end{itemize}
    
    \item \textbf{Selection Sort}: 
    \begin{itemize}
        \item \textbf{Time Complexity: }$O(n^2)$
        \item \textbf{Space Complexity: }Similar to bubble sort, selection sort is an in-place comparison-based sorting algorithm. It also requires only a constant amount of additional space For storing temporary variables. Hence, the space complexity of selection sort is \(O(1)\).
        \end{itemize}
    
    \item \textbf{Insertion Sort}: 
    \begin{itemize}
        \item \textbf{Time Complexity: }$O(n^2)$
        \item \textbf{Space Complexity: }Insertion sort is another in-place comparison-based sorting algorithm. Like bubble sort and selection sort, it requires only a constant amount of additional space For storing temporary variables. ThereFore, the space complexity of insertion sort is \(O(1)\).
        \end{itemize}
    
    \item \textbf{Merge Sort}:
    \begin{itemize}
        \item \textbf{Time Complexity: }$O(n \log n)$
        \item \textbf{Space Complexity: }Merge sort is a divide-and-conquer sorting algorithm that requires additional space For the merge operation. During the merging phase, it needs space to store the merged subarrays temporarily. The space complexity of merge sort is \(O(n)\), where \(n\) is the size of the input array.
        \end{itemize}
          
    
    \item \textbf{Quick Sort}: 
    \begin{itemize}
        \item \textbf{Time Complexity: }$O(n \log n)$ (average case), $O(n^2)$ (worst case)
        \item \textbf{Space Complexity: }Quick sort is another divide-and-conquer sorting algorithm that operates in-place. However, it typically requires additional space For the recursive function calls in the call stack. The worst-case space complexity of quick sort is \(O(n)\), where \(n\) is the size of the input array. In the average case, the space complexity is \(O(\log n)\).
        \end{itemize}
    \item \textbf{Heap Sort}:
    \begin{itemize}
        \item \textbf{Time Complexity: }$O(n \log n)$
        \item \textbf{Space Complexity: }Heap sort is an in-place comparison-based sorting algorithm that uses a binary heap data structure. It does not require any additional space beyond the input array size For storing temporary variables. ThereFore, the space complexity of heap sort is \(O(1)\).
        \end{itemize}
\end{itemize}
\FloatBarrier
\textbf{Choosing the Right Sorting Algorithm}

The choice of sorting algorithm depends on various factors such as the size of the input data, the distribution of data, memory constraints, and stability requirements. It is important to analyze these factors to select the most appropriate sorting algorithm For a given scenario.



\subsection{Graph Algorithms}
In computer engineering, graph algorithms are a set of techniques used to solve problems involving graphs. A graph is a collection of nodes (vertices) and edges that connect pairs of nodes. Graph algorithms can be broadly classIfied into two categories: traversal algorithms and path finding algorithms.

\subsubsection{Traversal Algorithms}

Traversal algorithms are used to visit or traverse all the nodes in a graph. The two most common traversal algorithms are depth-first search (DFS) and breadth-first search (BFS).
\begin{itemize}
    \item Depth-First Search (DFS):

Depth-first search explores as far as possible along each branch beFore backtracking. It starts at a selected node (the root) and explores as far as possible along each branch beFore backtracking.
\FloatBarrier
\begin{algorithm}
\caption{Depth-first Search (DFS)}
\begin{algorithmic}[1]
\Function{DFS}{$G, v, visited[]$}
    \State $visited[v] \gets \text{True}$
    \State \textbf{print} $v$
    \For{$u$ \textbf{in} $G.adjacentVertices(v)$}
        \If{$\neg visited[u]$}
            \State \Call{DFS}{$G, u, visited$}
        \EndIf
    \EndFor
\EndFunction
\end{algorithmic}
\end{algorithm}
\FloatBarrier
Python Code:
\FloatBarrier

\begin{verbatim}
def dfs(graph, v, visited):
    visited[v] = True
    print(v, end=' ')
    For u in graph.adjacent_vertices(v):
        If not visited[u]:
            dfs(graph, u, visited)
\end{verbatim}
\FloatBarrier

\item Breadth-First Search (BFS):

Breadth-first search explores all the neighboring nodes at the present depth beFore moving on to the nodes at the next depth level. It starts at a selected node (the root) and explores all of the neighbor nodes at the present depth beFore moving on to the nodes at the next depth level.
\FloatBarrier % Ensure algorithm appears after the paragraph
\begin{algorithm}
\caption{Breadth-first Search (BFS)}
\begin{algorithmic}[1]
\Function{BFS}{$G, start$}
    \State $visited[start] \gets \text{True}$
    \State $queue \gets [start]$
    \While{$queue$ is not empty}
        \State $v \gets$ dequeue from $queue$
        \State \textbf{print} $v$
        \For{$u$ \textbf{in} $G.adjacentVertices(v)$}
            \If{$\neg visited[u]$}
                \State $visited[u] \gets \text{True}$
                \State enqueue $u$ into $queue$
            \EndIf
        \EndFor
    \EndWhile
\EndFunction
\end{algorithmic}
\end{algorithm}
\FloatBarrier
Python Code:
\FloatBarrier
\begin{verbatim}
from collections import deque

def bfs(graph, start):
    visited = [False] * graph.num_vertices()
    visited[start] = True
    queue = deque([start])
    While queue:
        v = queue.popleft()
        print(v, end=' ')
        For u in graph.adjacent_vertices(v):
            If not visited[u]:
                visited[u] = True
                queue.append(u)
\end{verbatim}
\end{itemize}
\FloatBarrier % Ensure algorithm appears after the paragraph
\subsubsection{Pathfinding Algorithms}

Pathfinding algorithms are used to find the shortest path between two nodes in a graph. The most common pathfinding algorithm is Dijkstra's algorithm.
\FloatBarrier % Ensure algorithm appears after the paragraph
\begin{itemize}
\item Dijkstra's Algorithm:
Dijkstra's algorithm finds the shortest path between a starting node and all other nodes in the graph. 
\FloatBarrier % Ensure algorithm appears after the paragraph
\begin{algorithm}
\caption{Dijkstra's Algorithm}
\begin{algorithmic}[1]
\Function{Dijkstra}{$G, source$}
    \State $dist \gets$ dictionary of distances initialized to $\infty$ For all vertices
    \State $dist[source] \gets 0$
    \State $pq \gets$ priority queue initialized with $(0, source)$
    \While{$pq$ is not empty}
        \State $d, v \gets$ dequeue from $pq$
        \If{$d > dist[v]$}
            \State \textbf{continue}
        \EndIf
        \For{$u, w$ \textbf{in} $G.adjacentVertices(v)$}
            \State $new\_dist \gets dist[v] + w$
            \If{$new\_dist < dist[u]$}
                \State $dist[u] \gets new\_dist$
                \State enqueue $(new\_dist, u)$ into $pq$
            \EndIf
        \EndFor
    \EndWhile
\EndFunction
\end{algorithmic}
\end{algorithm}
\FloatBarrier
Python Code:
\FloatBarrier
\begin{verbatim}
import heapq

def dijkstra(graph, source):
    dist = {v: float('inf') For v in graph.vertices()}
    dist[source] = 0
    pq = [(0, source)]
    While pq:
        d, v = heapq.heappop(pq)
        If d > dist[v]:
            continue
        For u, w in graph.adjacent_vertices(v):
            new_dist = dist[v] + w
            If new_dist < dist[u]:
                dist[u] = new_dist
                heapq.heappush(pq, (new_dist, u))
\end{verbatim}
\end{itemize}
\FloatBarrier % Ensure algorithm appears after the paragraph

\section{Case Studies in Algorithm Analysis}
In algorithm analysis, three commonly used paradigms are \textit{Divide and Conquer Algorithms}, \textit{Dynamic Programming}, and \textit{Greedy Algorithms}.

\subsection{Divide and Conquer Algorithms}
Divide and Conquer Algorithms involve breaking down a problem into smaller subproblems, solving these subproblems recursively, and then combining the solutions to solve the original problem. The typical structure of a Divide and Conquer Algorithm includes three steps:
\begin{itemize}
  \item \textbf{Divide:} Break the problem into smaller subproblems.
  \item \textbf{Conquer:} Solve the subproblems recursively.
  \item \textbf{Combine:} Merge the solutions of the subproblems to solve the original problem.
\end{itemize}

\subsubsection{Analysis of Merge Sort and Quick Sort}

Merge sort and quick sort are two widely used sorting algorithms based on the divide-and-conquer paradigm. Both algorithms follow the typical structure of a divide-and-conquer algorithm, which involves three main steps: divide, conquer, and combine.

\paragraph{Merge Sort}

Merge sort is a stable sorting algorithm that divides the input array into two halves, sorts each half recursively, and then merges the sorted halves to produce a sorted output array. The key operation in merge sort is the merge step, where two sorted subarrays are merged into a single sorted array.
Algorithmic Example For Merge Sort:
     \FloatBarrier % Ensure algorithm appears after the paragraph
\begin{algorithm}
\caption{Merge Sort}
\begin{algorithmic}[1]
\Function{MergeSort}{$A[], start, end$}
    \If{$start < end$}
        \State $mid \gets (start + end) // 2$
        \State \Call{MergeSort}{$A[], start, mid$}
        \State \Call{MergeSort}{$A[], mid+1, end$}
        \State \Call{Merge}{$A[], start, mid, end$}
    \EndIf
\EndFunction
\FloatBarrier % Ensure algorithm appears after the paragraph
\Function{Merge}{$A[], start, mid, end$}
    \State $n_1 \gets mid - start + 1$
    \State $n_2 \gets end - mid$
    \State Create temporary arrays $L[1 \dots n_1]$ and $R[1 \dots n_2]$
    \For{$i \gets 1$ \textbf{to} $n_1$}
        \State $L[i] \gets A[start + i - 1]$
    \EndFor
    \For{$j \gets 1$ \textbf{to} $n_2$}
        \State $R[j] \gets A[mid + j]$
    \EndFor
    \State $i \gets 1$, $j \gets 1$, $k \gets start$
    \While{$i \leq n_1$ and $j \leq n_2$}
        \If{$L[i] \leq R[j]$}
            \State $A[k] \gets L[i]$
            \State $i \gets i + 1$
        \Else
            \State $A[k] \gets R[j]$
            \State $j \gets j + 1$
        \EndIf
        \State $k \gets k + 1$
    \EndWhile
    \While{$i \leq n_1$}
        \State $A[k] \gets L[i]$
        \State $i \gets i + 1$, $k \gets k + 1$
    \EndWhile
    \While{$j \leq n_2$}
        \State $A[k] \gets R[j]$
        \State $j \gets j + 1$, $k \gets k + 1$
    \EndWhile
\EndFunction
\end{algorithmic}
\end{algorithm}
\FloatBarrier % Ensure algorithm appears after the paragraph
Python Code:
\FloatBarrier % Ensure algorithm appears after the paragraph
\begin{verbatim}
def merge_sort(arr):
    If len(arr) > 1:
        mid = len(arr) // 2
        L = arr[:mid]
        R = arr[mid:]

        merge_sort(L)
        merge_sort(R)

        i = j = k = 0

        While i < len(L) and j < len(R):
            If L[i] < R[j]:
                arr[k] = L[i]
                i += 1
            else:
                arr[k] = R[j]
                j += 1
            k += 1

        While i < len(L):
            arr[k] = L[i]
            i += 1
            k += 1

        While j < len(R):
            arr[k] = R[j]
            j += 1
            k += 1
\end{verbatim}
    \FloatBarrier % Ensure algorithm appears after the paragraph

\paragraph{Time Complexity:}
The time complexity of merge sort is \(O(n \log n)\), where \(n\) is the size of the input array. This is because merge sort consistently divides the input array into halves until individual elements are reached, and then it merges these sorted halves. The merge operation takes \(O(n)\) time in the worst case.

\paragraph{Space Complexity:}
Merge sort typically requires additional space For the merge operation. The space complexity of merge sort is \(O(n)\), where \(n\) is the size of the input array, due to the need For temporary storage during the merge step.

\paragraph{Quick Sort}

Quick sort is an efficient in-place sorting algorithm that works by selecting a pivot element from the input array and partitioning the array into two subarrays, one containing elements less than the pivot and the other containing elements greater than the pivot. It then recursively sorts the subarrays.


Algorithmic Example For Quick Sort:
     \FloatBarrier % Ensure algorithm appears after the paragraph
    \begin{algorithm}
\caption{Quick Sort}
\begin{algorithmic}[1]
\Function{QuickSort}{$A[], start, end$}
    \If{$start < end$}
        \State $pivot \gets$ Partition($A[], start, end$)
        \State \Call{QuickSort}{$A[], start, pivot-1$}
        \State \Call{QuickSort}{$A[], pivot+1, end$}
    \EndIf
\EndFunction
\FloatBarrier % Ensure algorithm appears after the paragraph
\Function{Partition}{$A[], start, end$}
    \State $pivot \gets A[end]$
    \State $i \gets start - 1$
    \For{$j \gets start$ \textbf{to} $end-1$}
        \If{$A[j] < pivot$}
            \State $i \gets i + 1$
            \State Swap $A[i]$ and $A[j]$
        \EndIf
    \EndFor
    \State Swap $A[i + 1]$ and $A[end]$
    \State \Return $i + 1$
\EndFunction
\end{algorithmic}
\end{algorithm}
\FloatBarrier % Ensure algorithm appears after the paragraph
Python Code:
\begin{verbatim}
def quick_sort(arr, start, end):
    If start < end:
        pivot = partition(arr, start, end)
        quick_sort(arr, start, pivot - 1)
        quick_sort(arr, pivot + 1, end)

def partition(arr, start, end):
    pivot = arr[end]
    i = start - 1
    For j in range(start, end):
        If arr[j] < pivot:
            i += 1
            arr[i], arr[j] = arr[j], arr[i]
    arr[i + 1], arr[end] = arr[end], arr[i + 1]
    Return i + 1
\end{verbatim}
    \FloatBarrier % Ensure algorithm appears after the paragraph

\paragraph{Time Complexity:}
The time complexity of quick sort varies depending on the choice of pivot and the input array's initial order. In the average case, quick sort has a time complexity of \(O(n \log n)\). However, in the worst-case scenario, quick sort can degrade to \(O(n^2)\) If the pivot selection results in unbalanced partitions.

\paragraph{Space Complexity:}
Quick sort is an in-place sorting algorithm, meaning it does not require additional space beyond the input array size. ThereFore, the space complexity of quick sort is \(O(\log n)\) due to the recursive function calls in the call stack.

\paragraph{Comparison}

\paragraph{Time Complexity:}
In terms of time complexity, both merge sort and quick sort have an average-case time complexity of \(O(n \log n)\). However, quick sort may perForm worse in the worst-case scenario (\(O(n^2)\)) due to poor pivot selection, While merge sort consistently perForms in \(O(n \log n)\).

\paragraph{Space Complexity:}
Merge sort requires additional space For the merge operation, leading to a space complexity of \(O(n)\). In contrast, quick sort is an in-place algorithm and has a space complexity of \(O(\log n)\) due to the recursive calls.

In summary, both merge sort and quick sort are efficient sorting algorithms with similar average-case time complexities. Merge sort has better worst-case perFormance and stability but requires additional space, While quick sort is in-place but may perForm worse in certain scenarios.


\subsection{Dynamic Programming}
Dynamic Programming is a method For solving complex problems by breaking them down into simpler subproblems. It involves solving each subproblem only once and storing their solutions to avoid redundant computations. The key idea in Dynamic Programming is to use the solutions of overlapping subproblems.
\subsubsection{Analysis of Fibonacci Sequence Calculation}
A classic example of Dynamic Programming is the \textit{Fibonacci Sequence}. The Fibonacci Sequence can be efficiently computed using Dynamic Programming by storing the solutions to subproblems. The time complexity of the Dynamic Programming solution For the Fibonacci Sequence is $O(n)$.

The Fibonacci sequence is a series of numbers where each number is the sum of the two preceding ones, usually starting with 0 and 1. The sequence starts as follows: 0, 1, 1, 2, 3, 5, 8, 13, 21, ...

\paragraph{Algorithm}

The following algorithm uses dynamic programming to compute the $n^{th}$ Fibonacci number:
\FloatBarrier % Ensure algorithm appears after the paragraph
\begin{algorithm}
\caption{Fibonacci Sequence with Dynamic Programming}
\begin{algorithmic}[1]
\Procedure{Fibonacci}{$n$}
    \State Initialize an array $fib$ of size $n+1$
    \State $fib[0] \gets 0$, $fib[1] \gets 1$
    \For{$i \gets 2$ \textbf{to} $n$}
        \State $fib[i] \gets fib[i-1] + fib[i-2]$
    \EndFor
    \State \textbf{Return} $fib[n]$
\EndProcedure
\end{algorithmic}
\end{algorithm}
\FloatBarrier % Ensure algorithm appears after the paragraph
Python Code:
\FloatBarrier % Ensure algorithm appears after the paragraph
\begin{verbatim}
def fibonacci(n):
    fib = [0] * (n + 1)
    fib[0] = 0
    fib[1] = 1
    For i in range(2, n + 1):
        fib[i] = fib[i - 1] + fib[i - 2]
    Return fib[n]
\end{verbatim}
\FloatBarrier % Ensure algorithm appears after the paragraph
\paragraph{Explanation}

The algorithm computes the Fibonacci numbers iteratively using dynamic programming. It starts by initializing an array $fib$ of size $n+1$ to store the Fibonacci numbers.

The base cases are initialized as $fib[0] = 0$ and $fib[1] = 1$. Then, For each $i$ from $2$ to $n$, the algorithm calculates $fib[i]$ as the sum of the two preceding Fibonacci numbers, $fib[i-1]$ and $fib[i-2]$.

By iteratively computing and storing Fibonacci numbers, the algorithm avoids redundant calculations and improves efficiency.

\paragraph{Time Complexity}

The time complexity of the algorithm is $O(n)$, as it iterates through the array once to compute each Fibonacci number.
\paragraph{Space Complexity}
The space complexity of the given Fibonacci function can be analyzed as follows:

\begin{itemize}
    \item The function initializes a list `fib` of size `n + 1`, which requires space proportional to the input size `n`.
    \item ThereFore, the space complexity of the algorithm is \(O(n)\), as it uses additional space linearly dependent on the input size.
\end{itemize}

\subsection{Greedy Algorithms}
Greedy Algorithms make a series of choices based on a specIfic criterion at each step with the hope of finding an optimal solution. At each step, a greedy algorithm selects the best available option without considering future consequences. Greedy Algorithms do not always guarantee an optimal solution, but they often provide a good approximation.

An example of a Greedy Algorithm is the \textit{Fractional Knapsack Problem}. In this problem, items with weights and values are given, and the goal is to maximize the total value of items that can be fit into a knapsack of limited capacity. The algorithm selects items based on their value-to-weight ratio in a greedy manner.
In the Fractional Knapsack Problem, we are given a set of items, each with a weight $w_i$ and a value $v_i$, and a knapsack that can hold a maximum weight $W$. The objective is to maximize the total value of the items that can be placed into the knapsack without exceeding its weight capacity.

\subsubsection{Greedy Approach}

The greedy approach to solving the Fractional Knapsack Problem involves selecting items based on their value-to-weight ratio. At each step, we choose the item with the highest value-to-weight ratio and add as much of it as possible to the knapsack. This process is repeated until the knapsack is full or there are no more items left to consider.
\FloatBarrier % Ensure algorithm appears after the paragraph
\begin{algorithm}
\caption{Fractional Knapsack}
\begin{algorithmic}[1]
\Procedure{FractionalKnapsack}{$\text{items}, W$}
    \State Sort items by value-to-weight ratio in non-increasing order
    \State Initialize total value $V$ to $0$
    \State Initialize remaining capacity $C$ to $W$
    \For{each item $i$ in items}
        \If{$w_i \leq C$}
            \State Add item $i$ completely to the knapsack
            \State Update $V$ by adding $v_i$
            \State Update $C$ by subtracting $w_i$
        \Else
            \State Add a fraction of item $i$ to the knapsack
            \State Update $V$ by adding fraction of $v_i$
            \State Update $C$ to $0$
            \State Break
        \EndIf
    \EndFor
    \State \textbf{Return} $V$ \Comment{Total value of items in knapsack}
\EndProcedure
\end{algorithmic}
\end{algorithm}
\FloatBarrier % Ensure algorithm appears after the paragraph
Python Code:
\FloatBarrier % Ensure algorithm appears after the paragraph
\begin{verbatim}
def fractional_knapsack(items, W):
    items.sort(key=lambda x: x[1] / x[0], reverse=True)
    V = 0
    C = W
    For item in items:
        w_i, v_i = item
        If w_i <= C:
            V += v_i
            C -= w_i
        else:
            fraction = C / w_i
            V += fraction * v_i
            C = 0
            break
    Return V
\end{verbatim}
\FloatBarrier % Ensure algorithm appears after the paragraph
\paragraph{Mathematical Detail}

Let $n$ be the number of items.

\begin{itemize}
    \item The time complexity of the sorting step is $O(n \log n)$.
    \item The loop iterates over all $n$ items, each requiring $O(1)$ time For comparisons.
    \item ThereFore, the overall time complexity of the algorithm is $O(n \log n)$.
\end{itemize}

Let $V_{\text{opt}}$ be the optimal value that can be achieved.

The greedy algorithm may not always produce the optimal solution. However, it does guarantee a solution that is at least a fraction of the optimal solution. Let $V_{\text{greedy}}$ be the value obtained using the greedy algorithm.

\begin{itemize}
    \item Since the algorithm selects items based on their value-to-weight ratio, it ensures that each item contributes the maximum possible value For its weight.
    \item Thus, the value obtained by the greedy algorithm is at least $\frac{1}{2}$ times the optimal value, i.e., $V_{\text{greedy}} \geq \frac{1}{2} V_{\text{opt}}$.
\end{itemize}
\FloatBarrier % Ensure algorithm appears after the paragraph

\subsubsection{Kruskal's and Prim's Algorithms For Minimum Spanning Tree}

Kruskal's and Prim's algorithms are two popular greedy algorithms used to find the minimum spanning tree (MST) of a connected, undirected graph. 

\paragraph{Kruskal's Algorithm}

Kruskal's algorithm builds the MST by iteratively adding the smallest edge that does not Form a cycle until all vertices are connected.
\FloatBarrier
\begin{algorithm}
\caption{Kruskal's Algorithm}
\begin{algorithmic}[1]
\Function{Kruskal}{$G$}
    \State $T \gets \emptyset$ \Comment{Initialize empty MST}
    \State Sort edges of $G$ by weight
    \For{each edge $(u, v)$ in $G$ in sorted order}
        \If{$(u, v)$ does not create a cycle in $T$}
            \State Add $(u, v)$ to $T$
        \EndIf
    \EndFor
    \State \Return $T$
\EndFunction
\end{algorithmic}
\end{algorithm}
\textbf{Python Code:}
\FloatBarrier
\begin{verbatim}
def kruskal(G):
    T = []
    # Initialize empty MST
    # Sort edges of G by weight
    sorted_edges = sorted(G.edges(), key=lambda x: G[x[0]][x[1]]['weight'])
    For (u, v) in sorted_edges:
        # Check If adding (u, v) creates a cycle in T
        If not creates_cycle(T, (u, v)):
            # Add (u, v) to T
            T.append((u, v))
    Return T
\end{verbatim}
\FloatBarrier
\textbf{Example:} Consider the following graph $G$ with vertices $V = \{A, B, C, D, E\}$ and edges $E = \{(A, B, 5), (A, C, 3), (B, C, 1), (B, D, 4), (C, D, 2), (C, E, 6), (D, E, 7)\}$. Applying Kruskal's algorithm on $G$ results in the MST $T = \{(B, C, 1), (C, D, 2), (A, C, 3), (B, D, 4), (C, E, 6)\}$.

\paragraph{Prim's Algorithm}

Prim's algorithm grows the MST from a single vertex, iteratively adding the smallest edge connected to the current MST until all vertices are included.
\FloatBarrier
\begin{algorithm}
\caption{Prim's Algorithm}
\begin{algorithmic}[1]
\Function{Prim}{$G$}
    \State $T \gets \emptyset$ \Comment{Initialize empty MST}
    \State Choose a starting vertex $s$
    \State Initialize priority queue $Q$ with all vertices and keys $\infty$
    \State Set key of $s$ to 0
    \While{$Q$ is not empty}
        \State $u \gets$ Extract-Min($Q$)
        \For{each vertex $v$ adjacent to $u$}
            \If{$v$ is in $Q$ and weight of edge $(u, v)$ is less than key of $v$}
                \State Update key of $v$ to weight of edge $(u, v)$
                \State Update parent of $v$ to $u$
            \EndIf
        \EndFor
        \State Add edge $(u, \text{parent}[u])$ to $T$
    \EndWhile
    \State \Return $T$
\EndFunction
\end{algorithmic}
\end{algorithm}
\FloatBarrier
\textbf{Python Code:}
\FloatBarrier
\begin{algorithm}
\caption{Prim's Algorithm}
\begin{algorithmic}[1]
\Function{Prim}{$G$}
    \State $T \gets \emptyset$ \Comment{Initialize empty MST}
    \State Choose a starting vertex $s$
    \State Initialize priority queue $Q$ with all vertices and keys $\infty$
    \State Set key of $s$ to 0
    \While{$Q$ is not empty}
        \State $u \gets$ Extract-Min($Q$)
        \For{each vertex $v$ adjacent to $u$}
            \If{$v$ is in $Q$ and weight of edge $(u, v)$ is less than key of $v$}
                \State Update key of $v$ to weight of edge $(u, v)$
                \State Update parent of $v$ to $u$
            \EndIf
        \EndFor
        \State Add edge $(u, \text{parent}[u])$ to $T$
    \EndWhile
    \State \Return $T$
\EndFunction
\end{algorithmic}
\end{algorithm}
\FloatBarrier
\textbf{Example:} Using Prim's algorithm on the same graph $G$ yields the MST $T = \{(A, C, 3), (B, C, 1), (C, D, 2), (B, D, 4), (C, E, 6)\}$.



\section{Advanced Topics in Algorithmic Analysis}
\subsection{Randomized Algorithms:}

Randomized algorithms are a fascinating class of algorithms that introduce randomness into their execution process. Unlike deterministic algorithms, which follow a predefined sequence of steps to produce the same output For a given input, randomized algorithms incorporate randomness in their decision-making process. This randomness could be introduced through sources such as random number generators or probabilistic techniques.

One of the primary motivations For using randomized algorithms is their ability to provide efficient solutions to problems where deterministic algorithms may struggle due to their time complexity or complexity of design. By making random choices during execution, randomized algorithms can often achieve better average-case perFormance or simplIfy the problem-solving process.

The applications of randomized algorithms span a wide range of fields, including cryptography, machine learning, optimization problems, and computational biology, among others.
\FloatBarrier
\subsubsection{Randomized Quick Sort and Hashing}
\paragraph{Randomized QuickSort}

Randomized QuickSort is a randomized version of the QuickSort algorithm that randomly chooses a pivot element to partition the input array. This random choice helps in avoiding the worst-case time complexity of QuickSort.

Algorithm:
\FloatBarrier
\begin{algorithm}
\caption{Randomized QuickSort}
\begin{algorithmic}[1]
\State Array $A$, Indices $low$, $high$
\Function{randomizedQuickSort}{$A, low, high$}
    \If{$low < high$}
        \State $pivot \gets \text{random}(low, high)$
        \State $pivot \gets \text{partition}(A, low, high, pivot)$
        \State \textsc{randomizedQuickSort}($A, low, pivot-1$)
        \State \textsc{randomizedQuickSort}($A, pivot+1, high$)
    \EndIf
\EndFunction
\end{algorithmic}
\end{algorithm}

\FloatBarrier
\textbf{Python Code:}
\FloatBarrier
\begin{verbatim}
import random

def randomized_quick_sort(A, low, high):
    If low < high:
        pivot = random.randint(low, high)
        pivot = partition(A, low, high, pivot)
        randomized_quick_sort(A, low, pivot - 1)
        randomized_quick_sort(A, pivot + 1, high)

def partition(A, low, high, pivot):
    pivot_value = A[pivot]
    A[pivot], A[high] = A[high], A[pivot]
    i = low
    For j in range(low, high):
        If A[j] < pivot_value:
            A[i], A[j] = A[j], A[i]
            i += 1
    A[i], A[high] = A[high], A[i]
    Return i
\end{verbatim}
\FloatBarrier

Despite their advantages, randomized algorithms also present challenges, such as the need to analyze their probabilistic behavior and ensure correctness in all possible scenarios. Analyzing the perFormance of randomized algorithms often involves probabilistic reasoning and statistical techniques to estimate their expected behavior over a range of inputs.
\FloatBarrier
\paragraph{Hashing in Randomized Algorithms}

Hashing is a technique used to efficiently store and retrieve data in a data structure called a hash table. In the context of randomized algorithms, hashing involves mapping keys to indices in a hash table using a hash function. Randomized algorithms often use hashing For tasks such as maintaining sets, perForming dictionary lookups, or implementing caches.

\paragraph{Hash Function}

A hash function is a function that takes a key as input and produces an index (hash value) in the hash table. It should distribute keys unIFormly across the table to minimize collisions (two keys hashing to the same index). A good hash function has the following properties:

\begin{itemize}
    \item \textbf{Deterministic:} Given the same input key, the hash function should produce the same output index.
    \item \textbf{UnIFormity:} The hash function should evenly distribute keys across the hash table, reducing the likelihood of collisions.
    \item \textbf{Efficiency:} The hash function should be computationally efficient to calculate.
\end{itemize}

\paragraph{Collision Resolution}

Collisions occur when two keys hash to the same index in the hash table. There are various techniques to handle collisions, including:

\begin{itemize}
    \item \textbf{Chaining:} Each slot in the hash table contains a linked list of key-value pairs that hash to the same index.
    \item \textbf{Open Addressing:} In case of a collision, an alternative slot in the hash table is found using a probing sequence.
    \item \textbf{Double Hashing:} A second hash function is used to calculate the step size For probing, reducing clustering.
\end{itemize}

\paragraph{Algorithmic Example: Insertion with Chaining}

Let's consider an algorithm For inserting a key-value pair into a hash table using chaining For collision resolution:

\begin{algorithm}
\caption{Insertion with Chaining}
\begin{algorithmic}[1]
\Function{Insert}{$\text{hash\_table}, key, value$}
    \State $index \gets \text{hash\_function}(key)$
    \State Create a new node with key-value pair: $(key, value)$
    \State Append the node to the linked list at index $index$ in $\text{hash\_table}$
\EndFunction
\end{algorithmic}
\end{algorithm}

In this algorithm, the hash function maps the key to an index in the hash table. Then, a new node containing the key-value pair is inserted at the end of the linked list at the corresponding index.
\paragraph{Python Code: }
\begin{verbatim}
def Insert(hash_table, key, value):
    index = hash_function(key)
    new_node = Node(key, value)
    hash_table[index].append(new_node)
\end{verbatim}


\subsection{Parallel Algorithms:}

Parallel algorithms represent a pivotal paradigm in contemporary computing, harnessing the power of multiple processing units to enhance computational efficiency. As computing systems evolve with the prolIferation of multi-core processors and distributed environments, the signIficance of parallel algorithms becomes increasingly pronounced, facilitating concurrent execution of tasks and leveraging resources optimally.

One of the primary motivations behind parallel algorithms is to mitigate the inherent limitations of sequential processing, where tasks are executed sequentially, one after the other. By distributing the workload across multiple processing units, parallel algorithms enable tasks to be executed concurrently, thereby reducing overall execution time and improving system throughput.

Parallel algorithms find widespread applications across various domains, including scientIfic computing, data analytics, artIficial intelligence, and simulations. 
\subsubsection{Parallel Algorithm For Matrix Multiplication}
\FloatBarrier
Parallel matrix multiplication involves breaking down the matrix multiplication task into subtasks that can be computed in parallel across multiple processing units.
\FloatBarrier
\paragraph{Algorithm:}
\FloatBarrier

\begin{algorithm}
\caption{Parallel Matrix Multiplication}
\begin{algorithmic}
\Require{Matrices $A$ ($m \times n$), $B$ ($n \times p$)}

\Function{parallelMatrixMult}{$A, B$}
    \State $C \gets \text{initialize}(m \times p)$
    \For{$i$ from $1$ to $m$ in parallel}
        \For{$j$ from $1$ to $p$ in parallel}
            \For{$k$ from $1$ to $n$}
                \State $C[i][j] \mathrel{+}= A[i][k] \cdot B[k][j]$
            \EndFor
        \EndFor
    \EndFor
\EndFunction
\end{algorithmic}
\end{algorithm}
\FloatBarrier
\paragraph{Python Code:}
\FloatBarrier
\begin{verbatim}
import multiprocessing

def parallel_matrix_mult(A, B):
    m, n = len(A), len(A[0])
    p = len(B[0])
    C = [[0] * p For _ in range(m)]

    def compute_cell(i, j):
        cell_value = 0
        For k in range(n):
            cell_value += A[i][k] * B[k][j]
        Return cell_value

    def update_cell(i, j, cell_value):
        C[i][j] = cell_value

    def compute_row(i):
        For j in range(p):
            cell_value = compute_cell(i, j)
            update_cell(i, j, cell_value)

    num_cores = multiprocessing.cpu_count()
    pool = multiprocessing.Pool(processes=num_cores)
    pool.map(compute_row, range(m))
    pool.close()
    pool.join()

    Return C
\end{verbatim}
\FloatBarrier
Designing efficient parallel algorithms involves addressing several challenges, including load balancing, synchronization, communication overhead, and scalability. Effective parallelization requires careful consideration of algorithmic design, data partitioning strategies, and parallel execution models to achieve optimal perFormance across diverse computing architectures.

In conclusion, parallel algorithms represent a cornerstone of modern computing, enabling efficient utilization of computational resources and scalability across diverse application domains. As computing systems continue to evolve with the advent of multi-core processors, accelerators, and distributed architectures, the importance of parallel algorithms in driving computational advancements and addressing complex challenges will continue to grow.


\subsection{Approximation Algorithms:}

Approximation algorithms represent a powerful approach to solving optimization problems when finding an exact solution is prohibitively expensive or impractical. These algorithms trade off computational complexity For efficiency by providing solutions that are near-optimal within a guaranteed approximation factor. By sacrIficing optimality For speed, approximation algorithms enable the rapid solution of complex optimization problems in various fields, including computer science, operations research, and engineering.

The versatility of approximation algorithms makes them indispensable in scenarios where exact solutions are either unattainable or require excessive computational resources.

\subsubsection{Travelling Salesman Problem and Vertex Cover}
\paragraph{Travelling Salesman Problem}
The Travelling Salesman Problem (TSP) is a classic optimization problem in which the goal is to find the shortest possible route that visits each city exactly once and Returns to the original city. It is a well-known NP-hard problem, meaning that there is no known polynomial-time algorithm that can solve it optimally For all instances. However, there exist approximation algorithms that provide near-optimal solutions within a guaranteed approximation factor.

One such approximation algorithm For the TSP is the Nearest Neighbor Algorithm. The algorithm starts at a random city and repeatedly selects the nearest unvisited city until all cities have been visited. Finally, it Returns to the starting city, completing the tour. Although the Nearest Neighbor Algorithm does not guarantee the optimal solution, it often produces solutions that are relatively close to the optimal.

\textbf{Algorithmic Example: Nearest Neighbor Algorithm For TSP}

\begin{algorithm}
\caption{Nearest Neighbor Algorithm For TSP}
\begin{algorithmic}[1]
\Function{NearestNeighborTSP}{$G, start$}
    \State $T \gets \emptyset$ \Comment{Initialize empty tour}
    \State $current \gets start$ \Comment{Start from the given city}
    \While{there are unvisited cities}
        \State Add $current$ to $T$ \Comment{Visit the current city}
        \State $next \gets$ nearest unvisited city to $current$ \Comment{Find the nearest unvisited city}
        \State $current \gets next$ \Comment{Move to the nearest city}
    \EndWhile
    \State Add edge from last city to $start$ to complete the tour
    \State \Return $T$
\EndFunction
\end{algorithmic}
\end{algorithm}
\textbf{Python Code:}
\FloatBarrier
\begin{verbatim}
def nearest_neighbor_tsp(G, start):
    T = []  # Initialize empty tour
    current = start  # Start from the given city
    While there are unvisited cities:
        T.append(current)  # Visit the current city
        next = nearest_unvisited_city_to(current)  # Find the nearest unvisited city
        current = next  # Move to the nearest city
    T.append(edge_from_last_city_to_start)  # Add edge from last city to start to complete the tour
    Return T
\end{verbatim}
\FloatBarrier
\paragraph{Vertex Cover Problem}
The Vertex Cover problem is another well-known optimization problem where the objective is to find the smallest set of vertices that covers all edges in a graph. Similar to the TSP, the Vertex Cover problem is NP-hard, and finding an exact solution is computationally expensive. However, various approximation algorithms exist For Vertex Cover that provide solutions within a certain approximation factor.

One such approximation algorithm For Vertex Cover is the Greedy Algorithm. The Greedy Algorithm iteratively selects vertices that cover the maximum number of uncovered edges until all edges are covered. While the Greedy Algorithm may not yield the optimal solution, it guarantees a solution within twice the size of the optimal Vertex Cover.

\textbf{Algorithmic Example: Greedy Algorithm For Vertex Cover}

\begin{algorithm}
\caption{Greedy Algorithm For Vertex Cover}
\begin{algorithmic}[1]
\Function{GreedyVertexCover}{$G$}
    \State $C \gets \emptyset$ \Comment{Initialize empty vertex cover}
    \State $E' \gets$ all edges of $G$ \Comment{Initialize uncovered edges}
    \While{$E'$ is not empty}
        \State Select an arbitrary edge $e = (u, v)$ from $E'$
        \State Add vertices $u$ and $v$ to $C$
        \State Remove all edges incident to $u$ and $v$ from $E'$
    \EndWhile
    \State \Return $C$
\EndFunction
\end{algorithmic}
\end{algorithm}
\textbf{Python Code:}
\FloatBarrier
\begin{verbatim}
def greedy_vertex_cover(G):
    C = set()  # Initialize empty vertex cover
    E_prime = G.edges()  # Initialize uncovered edges
    While E_prime:  # While there are uncovered edges
        u, v = E_prime.pop()  # Select an arbitrary edge
        C.add(u)  # Add vertices u and v to the vertex cover
        C.add(v)
        # Remove all edges incident to u and v from the set of uncovered edges
        E_prime.dIfference_update(G.edges(u))
        E_prime.dIfference_update(G.edges(v))
    Return C
\end{verbatim}
\FloatBarrier
These approximation algorithms For the Travelling Salesman Problem and Vertex Cover demonstrate how approximation techniques can be used to efficiently find near-optimal solutions to NP-hard optimization problems.


\section{Tools and Techniques For Algorithm Analysis}
\subsection{Profiling and Benchmarking Tools}

Profiling and benchmarking tools are crucial For delving into the inner workings of algorithms and assessing their perFormance. Profiling tools allow developers to pinpoint inefficiencies and bottlenecks within their code, While benchmarking tools provide a means to compare the perFormance of dIfferent algorithms or implementations against specIfied metrics.

\subsubsection{Profiling Tools}
Profiling tools gather and analyze data regarding the runtime behavior of a program. This data typically includes details such as execution time of various code segments, memory utilization, and patterns of function calls. By utilizing profiling tools, developers can identIfy areas For optimization and enhance the overall perFormance of their algorithms.
\FloatBarrier
Example: 
\FloatBarrier
One prominent profiling tool is `cProfile` in Python, which furnishes detailed insights into the function call stack and execution time of individual functions.
\FloatBarrier
\begin{verbatim}
import cProfile

def example_function():
    For i in range(1000000):
        pass

cProfile.run('example_function()')
\end{verbatim}
\FloatBarrier
\subsubsection{Benchmarking Tools}
Benchmarking tools are invaluable For comparing and evaluating the perFormance of diverse algorithms or implementations. These tools facilitate the measurement of metrics such as execution time, memory consumption, and scalability, enabling developers to make inFormed decisions based on perFormance analyses.
\FloatBarrier
Example: 
\FloatBarrier
`pytest-benchmark`, a Python package, offers functionality For benchmarking functions and facilitating perFormance comparisons.
\FloatBarrier
\begin{verbatim}
import pytest

def test_perFormance(benchmark):
    result = benchmark(some_function, arg1, arg2)

    assert result == expected_output
\end{verbatim}
\FloatBarrier
\subsection{Simulation and Modeling Tools}

Simulation and modeling tools play a pivotal role in algorithm analysis by enabling the creation of simulated environments For evaluating algorithm behavior under varying conditions.

\subsubsection{Simulation Tools}
Simulation tools empower developers to simulate real-world scenarios within controlled virtual environments. These tools aid in testing algorithms across a spectrum of scenarios to assess their robustness and efficiency.
\FloatBarrier
Example: 
\FloatBarrier
`SimPy`, a process-based discrete-event simulation framework in Python, is widely utilized For simulating diverse scenarios.
\FloatBarrier
\begin{verbatim}
import simpy

def example_simulation(env):
    resource = simpy.Resource(env, capacity=1)
    with resource.request() as request:
        yield request
        # Simulation logic here
\end{verbatim}
\FloatBarrier

\subsubsection{Modeling Tools}
Modeling tools are instrumental in representing algorithms, systems, or processes through mathematical models For analysis and optimization purposes. These tools assist in creating abstract representations that can be analyzed For perFormance assessment and improvement.
\FloatBarrier
Example:
\FloatBarrier
`MATLAB` is a prevalent tool renowned For modeling algorithms and systems through mathematical equations and simulations.
\FloatBarrier
\begin{verbatim}
function output = example_model(input)
    % Model equations and logic here
end
\end{verbatim}
\FloatBarrier
These diverse tools and techniques serve as invaluable assets For the analysis, optimization, and improvement of algorithms, providing developers with the means to enhance the efficiency and perFormance of their computational solutions.

\section{Conclusion}
In this section, we will discuss the importance of algorithm analysis and highlight some future trends in algorithm analysis. Understanding the efficiency of algorithms is crucial For designing scalable and optimal solutions to computational problems. Additionally, staying abreast of the latest advancements in algorithm analysis is essential For keeping pace with the rapidly evolving field of computer science.

\subsection{The Importance of Algorithm Analysis}
Algorithm analysis plays a fundamental role in computer science as it enables us to compare dIfferent algorithms based on their efficiency and perFormance. One of the key aspects of algorithm analysis is time complexity, which quantIfies the amount of time an algorithm takes to run as a function of the input size. For example, consider the comparison between linear search and binary search algorithms. Linear search has a time complexity of O(n) where n is the size of the input array, While binary search has a time complexity of O(log n). By analyzing these complexities, we can understand the perFormance trade-offs between these two algorithms and choose the most efficient one For a given problem.

\subsection{Future Trends in Algorithm Analysis}
The field of algorithm analysis is constantly evolving, with researchers exploring new techniques to analyze and optimize algorithms. One emerging trend is the use of machine learning algorithms to predict the perFormance of dIfferent algorithms on specIfic inputs. By training models on historical data, researchers can develop predictive algorithms that recommend the most suitable algorithm For a given problem instance. Another trend is the application of quantum computing principles to algorithm analysis. Quantum algorithms offer the potential For exponential speedup over classical algorithms in certain problem domains, opening up new possibilities For solving complex computational problems efficiently.


\section{Further Reading and Resources}
In the field of algorithm analysis, a wealth of resources is available to deepen your understanding and refine your skills. These resources range from foundational textbooks and seminal research papers to comprehensive online tutorials and courses. Additionally, open-source libraries provide practical tools and implementations to aid in the study and application of algorithms. The following subsections provide a detailed guide to these valuable resources.

\subsection{Key Papers and Books on Algorithms}
To gain a solid foundation in algorithm analysis, it is crucial to refer to key papers and authoritative textbooks that have shaped the field. 

One of the cornerstone texts is "Introduction to Algorithms" by Cormen, Leiserson, Rivest, and Stein. This book, often referred to as CLRS, covers a wide range of algorithms and provides in-depth analysis techniques . Another essential book is "The Algorithm Design Manual" by Steven S. Skiena, which offers practical guidance on designing and implementing algorithms, along with a comprehensive catalog of algorithmic problems and solutions .

For more advanced topics, "Algorithms" by Robert Sedgewick and Kevin Wayne provides a deep dive into data structures, sorting, searching, and graph algorithms, complemented by practical implementations in Java . Additionally, "Algorithmic Graph Theory and Perfect Graphs" by Martin Charles Golumbic explores the intersection of graph theory and algorithm design .

Seminal papers such as "Computers and Intractability: A Guide to the Theory of NP-Completeness" by Michael Garey and David Johnson provide a foundational understanding of computational complexity and NP-completeness, which are crucial For algorithm analysis .

\subsection{Online Tutorials and Courses}
With the prolIferation of online learning platForms, numerous high-quality tutorials and courses are available to help you master algorithm analysis. 

Coursera offers several courses taught by experts from leading universities. The "Algorithms Specialization" by StanFord University, taught by Tim Roughgarden, covers essential algorithmic techniques and their applications . Similarly, the "Algorithmic Toolbox" by the University of CalIFornia, San Diego, and National Research University Higher School of Economics provides a practical approach to algorithm design and analysis .

On edX, the Massachusetts Institute of Technology (MIT) offers the "Introduction to Computer Science and Programming Using Python" course, which includes a module on algorithms and their complexity . Additionally, Khan Academy provides free tutorials on various algorithmic concepts, from basic sorting and searching to dynamic programming and graph algorithms .

PlatForms like Udacity also offer nanodegree programs in data structures and algorithms, providing hands-on projects and mentorship to help solidIfy your understanding and skills .

\subsection{Open-Source Libraries and Algorithm Implementations}
Open-source libraries are invaluable resources For studying and implementing algorithms. These libraries offer pre-implemented algorithms, allowing you to focus on understanding and applying them to solve problems.

\begin{itemize}
    \item \textbf{NetworkX}: A Python library For the creation, manipulation, and study of complex networks of nodes and edges. It includes implementations of many graph algorithms.
    \item \textbf{Scikit-learn}: A machine learning library in Python that includes tools For data mining and data analysis, built on NumPy, SciPy, and matplotlib. It provides implementations of many standard machine learning algorithms.
    \item \textbf{Boost C++ Libraries}: A collection of peer-reviewed, open-source C++ libraries that provide implementations of various algorithms and data structures.
    \item \textbf{Graph-tool}: An efficient Python library For manipulation and statistical analysis of graphs, which includes a variety of graph algorithms.
\end{itemize}

Let's take a detailed look at how some of these libraries can be used to implement and analyze algorithms. Consider the example of using NetworkX to implement and analyze Dijkstra's algorithm For finding the shortest path in a graph.

\FloatBarrier
\begin{lstlisting}[language=Python, caption=Using NetworkX to Implement Dijkstra's Algorithm]
import networkx as nx

# Create a directed graph
G = nx.DiGraph()

# Add edges along with their weights
G.add_weighted_edges_from([
    ('A', 'B', 1),
    ('B', 'C', 2),
    ('A', 'C', 4),
    ('C', 'D', 1),
    ('D', 'E', 5)
])

# Compute shortest path from node 'A' to all other nodes
shortest_paths = nx.single_source_dijkstra_path(G, source='A')

# Print the shortest paths
print(shortest_paths)
\end{lstlisting}
\FloatBarrier

In this example, we use the NetworkX library to create a directed graph and add weighted edges to it. We then use the built-in function 'single\_source\_dijkstra\_path' to compute the shortest path from node 'A' to all other nodes. This demonstrates how open-source libraries can simplIfy the implementation and analysis of complex algorithms.

\section{End of Chapter Exercises}
In this section, we provide a series of exercises designed to reinForce and extend the concepts covered in this chapter on Algorithm Analysis. These exercises are divided into two main categories: conceptual questions to reinForce learning and practical problems with solutions. The goal of these exercises is to deepen your understanding of algorithm analysis through both theoretical questions and hands-on coding problems.

\subsection{Conceptual Questions to ReinForce Learning}
This subsection presents a series of conceptual questions aimed at reinForcing the fundamental concepts discussed in this chapter. These questions are designed to test your comprehension and encourage you to think critically about the material.

\begin{itemize}
    \item \textbf{What is the time complexity of the following algorithm, and how do you derive it?} 
    \begin{algorithm}
    \caption{Sample Algorithm}
    \begin{algorithmic}
        \For{$i = 1$ to $n$}
            \For{$j = 1$ to $n$}
                \State PerForm constant time operation
            \EndFor
        \EndFor
    \end{algorithmic}
    \end{algorithm}
    \item \textbf{Explain the dIfference between Big-O, Big-Theta, and Big-Omega notation. Provide examples.}
    \item \textbf{How does divide-and-conquer strategy improve the efficiency of algorithms? Illustrate with the example of Merge Sort.}
    \item \textbf{Describe the Master Theorem For solving recurrences. Give an example of how it is applied.}
    \item \textbf{What are the trade-offs between time complexity and space complexity? Provide examples of algorithms where these trade-offs are evident.}
    \item \textbf{Define amortized analysis and explain how it dIffers from average-case analysis. Provide an example using the dynamic array.}
\end{itemize}

\subsection{Practical Problems and Solutions}
This subsection provides practical problems related to algorithm analysis, along with detailed solutions. These problems will require you to apply the concepts learned in the chapter and implement algorithms in Python. Each problem is followed by a solution that includes both the algorithmic description and the Python code.

\begin{itemize}
    \item \textbf{Problem 1: Implement and analyze the time complexity of Binary Search.}
    \begin{itemize}
        \item \textbf{Description:} Write a Python function that implements binary search. Given a sorted array and a target value, the function should Return the index of the target value If it is present in the array, or -1 If it is not.
        \item \textbf{Algorithm:}
        \begin{algorithm}
        \caption{Binary Search}
        \begin{algorithmic}
            \State Let $low \leftarrow 0$
            \State Let $high \leftarrow n-1$
            \While{$low \leq high$}
                \State $mid \leftarrow \left\lfloor \frac{low + high}{2} \right\rfloor$
                \If{$array[mid] == target$}
                    \Return $mid$
                \ElsIf{$array[mid] < target$}
                    \State $low \leftarrow mid + 1$
                \Else
                    \State $high \leftarrow mid - 1$
                \EndIf
            \EndWhile
            \Return -1
        \end{algorithmic}
        \end{algorithm}
        \item \textbf{Python Code:}
        \FloatBarrier
        \begin{lstlisting}[language=Python]
def binary_search(array, target):
    low = 0
    high = len(array) - 1

    While low <= high:
        mid = (low + high) // 2
        If array[mid] == target:
            Return mid
        elIf array[mid] < target:
            low = mid + 1
        else:
            high = mid - 1

    Return -1
        \end{lstlisting}
        \FloatBarrier
        \item \textbf{Time Complexity:} The time complexity of binary search is $O(\log n)$, where $n$ is the number of elements in the array. This is because the algorithm repeatedly divides the search interval in half.
    \end{itemize}

    \item \textbf{Problem 2: Analyze the perFormance of a sorting algorithm (e.g., Quick Sort).}
    \begin{itemize}
        \item \textbf{Description:} Implement Quick Sort in Python and analyze its average and worst-case time complexities.
        \item \textbf{Algorithm:}
        \begin{algorithm}
        \caption{Quick Sort}
        \begin{algorithmic}
            \State \textbf{Function} \textsc{QuickSort}(array, low, high)
            \If{$low < high$}
                \State $pi \leftarrow$ \textsc{Partition}(array, low, high)
                \State \textsc{QuickSort}(array, low, pi - 1)
                \State \textsc{QuickSort}(array, pi + 1, high)
            \EndIf
        \end{algorithmic}
        \end{algorithm}
        \begin{algorithm}
        \caption{Partition}
        \begin{algorithmic}
            \State \textbf{Function} \textsc{Partition}(array, low, high)
            \State $pivot \leftarrow array[high]$
            \State $i \leftarrow low - 1$
            \For{$j \leftarrow low$ to $high - 1$}
                \If{$array[j] < pivot$}
                    \State $i \leftarrow i + 1$
                    \State Swap $array[i]$ and $array[j]$
                \EndIf
            \EndFor
            \State Swap $array[i + 1]$ and $array[high]$
            \Return $i + 1$
        \end{algorithmic}
        \end{algorithm}
        \item \textbf{Python Code:}
        \FloatBarrier
        \begin{lstlisting}[language=Python]
def partition(array, low, high):
    pivot = array[high]
    i = low - 1
    For j in range(low, high):
        If array[j] < pivot:
            i += 1
            array[i], array[j] = array[j], array[i]
    array[i + 1], array[high] = array[high], array[i + 1]
    Return i + 1

def quick_sort(array, low, high):
    If low < high:
        pi = partition(array, low, high)
        quick_sort(array, low, pi - 1)
        quick_sort(array, pi + 1, high)

# Example usage
array = [10, 7, 8, 9, 1, 5]
n = len(array)
quick_sort(array, 0, n-1)
print("Sorted array:", array)
        \end{lstlisting}
        \FloatBarrier
        \item \textbf{Time Complexity:}
        \begin{itemize}
            \item \textbf{Average Case:} $O(n \log n)$
            \item \textbf{Worst Case:} $O(n^2)$, occurs when the smallest or largest element is always chosen as the pivot.
        \end{itemize}
    \end{itemize}

    \item \textbf{Problem 3: Implement and analyze the time complexity of Dijkstra's Algorithm.}
    \begin{itemize}
        \item \textbf{Description:} Write a Python function to find the shortest path from a source node to all other nodes in a weighted graph using Dijkstra's algorithm.
        \item \textbf{Algorithm:}
        \begin{algorithm}
        \caption{Dijkstra's Algorithm}
        \begin{algorithmic}
            \State \textbf{Function} \textsc{Dijkstra}(graph, source)
            \State Initialize distance of all vertices as infinity
            \State Distance of source vertex from itself is always 0
            \State Create a priority queue and add source to it
            \While{priority queue is not empty}
                \State Extract vertex $u$ with minimum distance
                \For{each neighbor $v$ of $u$}
                    \If{distance to $v$ through $u$ is less than current distance of $v$}
                        \State Update distance of $v$
                        \State Add $v$ to the priority queue
                    \EndIf
                \EndFor
            \EndWhile
            \Return distances from source to all vertices
        \end{algorithmic}
        \end{algorithm}
        \item \textbf{Python Code:}
        \FloatBarrier
        \begin{lstlisting}[language=Python]
import heapq

def dijkstra(graph, start):
    priority_queue = []
    heapq.heappush(priority_queue, (0, start))
    distances = {vertex: float('infinity') For vertex in graph}
    distances[start] = 0

    While priority_queue:
        current_distance, current_vertex = heapq.heappop(priority_queue)

        If current_distance > distances[current_vertex]:
            continue

        For neighbor, weight in graph[current_vertex].items():
            distance = current_distance + weight

            If distance < distances[neighbor]:
                distances[neighbor] = distance
                heapq.heappush(priority_queue, (distance, neighbor))

    Return distances

# Example usage
graph = {
    'A': {'B': 1, 'C': 4},
    'B': {'A': 1, 'C': 2, 'D': 5},
    'C': {'A': 4, 'B': 2, 'D': 1},
    'D': {'B': 5, 'C': 1}
}
print(dijkstra(graph, 'A'))
        \end{lstlisting}
        \FloatBarrier
        \item \textbf{Time Complexity:} Using a priority queue, the time complexity is $O((V + E) \log V)$ where $V$ is the number of vertices and $E$ is the number of edges.
    \end{itemize}
\end{itemize}


%\section*{References}
\begin{thebibliography}{9}
\bibitem{knuth}
Knuth: Recent News. \textit{Donald Knuth}. [Online]. Available: \url{http://www-cs-faculty.stanFord.edu/~uno/news.html}. [Accessed: 28 August 2016].
\bibitem{AhoHopcroft1974}
Alfred V. Aho, John E. Hopcroft, Jeffrey D. Ullman. \textit{The design and analysis of computer algorithms}. Addison-Wesley Pub. Co., 1974.
\bibitem{Hromkovič2004}
Juraj Hromkovič. \textit{Theoretical computer science: introduction to Automata, computability, complexity, algorithmics, randomization, communication, and cryptography}. Springer, 2004.
\bibitem{Ausiello1999}
Giorgio Ausiello. \textit{Complexity and approximation: combinatorial optimization problems and their approximability properties}. Springer, 1999.
\bibitem{Wegener20}
Ingo Wegener. \textit{Complexity theory: exploring the limits of efficient algorithms}. Springer-Verlag, 2005.
\bibitem{Tarjan1983}
Robert Endre Tarjan. \textit{Data structures and network algorithms}. SIAM, 1983.
\bibitem{PriceAbstraction}
"Examples of the price of abstraction?". \textit{cstheory.stackexchange.com}. [Online]. Available: \url{https://cstheory.stackexchange.com/q/608}.
\bibitem{PowerRule}
R. J. Lipton. "How To Avoid O-Abuse and Bribes". [Online]. Available: \url{http://rjlipton.wordpress.com/2009/07/24/how-to-avoid-o-abuse-and-bribes/}.
\end{thebibliography}

\end{document}
